% =====================================================================
%  R. Nielsen, ON THE POWER OF IMAGINATION.  Six Lectures.  Copenhagen: Gyldendal (F. Hegel), 1876.  91 pp.
%  (Om Phantasiens Magt.  Sex Forelæsninger.)
%  ENGLISH TRANSLATION, from transcription.tex (Danish).
%  GENERATED by ebook-method/phantasiens-magt/tr_build.py --write from this template and the four translators' parts
%  in ebook-method/phantasiens-magt/tr/ (TRANSLATOR-BRIEF.md; terms: GLOSSARY.md, Phantasie = imagination and the title
%  approved by Hans 2026-10-12).  Conventions: TRANSLATION-PLAYBOOK.md.  Page numbers of the 1876 print in the margin.
% =====================================================================
\documentclass[12pt,a4paper,openany]{book}
\usepackage[utf8]{inputenc}
\usepackage[T1]{fontenc}
\usepackage{libertinus}
\usepackage{libertinust1math}
\usepackage[english]{babel}
\usepackage[protrusion=true,expansion=true]{microtype}
\usepackage{setspace}
\usepackage[margin=1.3in]{geometry}
\usepackage{fancyhdr}
\usepackage{hyperref}
\hypersetup{pdftitle={On the Power of Imagination}, pdfauthor={Rasmus Nielsen}, hidelinks}
\setlength{\headheight}{15pt}
\pagestyle{fancy}
\fancyhf{}
\fancyhead[LE,RO]{\thepage}
\fancyhead[C]{\textit{R. Nielsen: On the Power of Imagination (1876)}}
\setstretch{1.3}
\usepackage{marginnote}
\newcommand{\opage}[1]{\marginnote{\footnotesize\textit{[#1]}}}
% a lecture opens a new page: heading, a short rule
\newcommand{\lecture}[1]{\clearpage\begin{center}{\large #1}\\[0.4em]\rule{3em}{0.4pt}\end{center}\medskip\noindent}

\begin{document}

% PDF 6, the title page
\begin{titlepage}
\begin{center}
\vspace*{3em}
{\large ON}\\[1em]
{\Huge THE POWER OF IMAGINATION.}\\[2.5em]
SIX LECTURES\\[1em]
{\small BY}\\[1em]
{\large R. NIELSEN.}\\[5em]
COPENHAGEN.\\[0.5em]
PUBLISHED BY GYLDENDAL'S BOOKSHOP (F. HEGEL).\\[0.5em]
{\small PRINTED BY J. JØRGENSEN \& CO.}\\[0.5em]
{\small 1876.}\\[3em]
{\small Translated from the Danish}
\end{center}
\end{titlepage}

% PDF 8, the contents
\begin{center}{\large Contents.}\\[0.3em]\rule{3em}{0.4pt}\end{center}
\hfill{\small Page}\par
\noindent\emph{First Lecture.}\enspace Imagination and Reality\dotfill\ 1.\par
\noindent\emph{Second Lecture.}\enspace The Relation of Imagination to the Question\par
\noindent\hspace*{2em}of the Immortality of the Soul\dotfill\ 14.\par
\noindent\emph{Third Lecture.}\enspace Religious Imagination\dotfill\ 28.\par
\noindent\emph{Fourth Lecture.}\enspace Old Testament Imagination\dotfill\ 43.\par
\noindent\emph{Fifth Lecture.}\enspace New Testament Imagination\dotfill\ 57.\par
\noindent\emph{Sixth Lecture.}\enspace The Infinite Realm of Imagination\dotfill\ 75.\par
\clearpage

% ===== TEXT (pp. 1--91) =====
% --- p. 1 ---
\setcounter{footnote}{0}%
\lecture{First Lecture.}\opage{1}

We know daily life, its joys and sorrows, its comfort and peace. Everything here is so homelike to us, so accustomed,
so familiar. But behind the familiar there is an unknown, and this Unknown sometimes announces itself involuntarily
with questions that cause anxiety. We sit in the family circle in the sunshine of domestic happiness; the Unknown
whispers: how long will it last? We remember our departed ones, and religion comforts us with a reunion; the Unknown
whispers: is there a reunion, are you certain of it? So in manifold cases. Now what is this Unknown? Is it of evil or
of good? Whence does it stem, from the source of light or from the regions of darkness? It comes as an express
messenger, an omen, a reminder of the invisible, an admonishing voice from the world of spirits. There are people who
have a natural shyness of making closer acquaintance with this Unknown. But everyday life itself is subject to
dispensations that can bring the Unknown to step forth, suddenly and in all its might. Let me cite a little example.

A young man, with original gifts and destined for studies, had already gone through the
% --- p. 2 ---
\setcounter{footnote}{0}%
\opage{2}first academic course, when his father, who was beginning to feel old and tired, wished that his son should
become a farmer and help him manage the farm. The son gave in to his father's wish, all the more willingly as there was
a manor in the district, and on the manor a young lady, the landowner's daughter. The young man had looked at her, and
she in turn had looked at him --- his decision was taken. Toward autumn --- he had by that time already been a couple
of years in his new occupation --- a comet appeared in the sky. It gave people something to think about and set the
minds of the country folk astir. Old women assured everyone that war, famine, earthquake and pestilence were impending,
either all four things or at least one of them; it could never fail. The younger ones thought that the matter was surely
not so dangerous; the comet had, after all, been duly announced in the almanac. They read and saw that week after week
it punctually and quietly conformed to the almanac's prescriptions: it could not possibly have anything evil in mind.
But as it came nearer week by week, the wonder grew, for the comet seemed to change and assume a threatening
appearance. The red eye grew ever sharper, the tail grew, spread and branched like a great rod of fire. The stars in
the clear September sky took on a duller light; it was as if heaven wanted half to close its eyes so as not to see the
uncanny sign. The old women were seized by a prophetic spirit; they preached penance and conversion, and many became
devout.

Yet natural life asserts its rights. The harvest had been blessed; of famine there was thus no danger. The rowan
berries turned red, the foliage in the wood turned brown and violet and yellow, and the leaves
% --- p. 3 ---
\setcounter{footnote}{0}%
\opage{3}fell; the wind did indeed rush breathlessly over stubble and ploughed fields, but of earthquake there was
nothing to be noticed; the state of health was good, and the newspapers could surely vouch for the peace: so there was
to be a harvest feast after all.

The harvest feast at the manor was the neighborhood's annual and most brilliant festival; it was in truth a popular
festival, for the liberal landowner invited high and low. Maidens and matrons, noblemen and peasants mingled with one
another in the dance; the stream surged in the great halls and through the garden lit with torches and lamps. Many
had looked forward to this festival, but none more than our young friend. He now expected so surely that this evening
was to crown his happiness. As he rode along the road he formed a resolution that was to stand firm; he would this
evening declare his secret engagement to the landowner; he had carefully chosen his words and already saw in spirit
how the distinguished but humane gentleman, who had always appreciated his enterprise and shown him so much kindness,
let the first wrinkles of surprise smooth themselves out and gave in to the united entreaties of himself and the
beloved daughter. It seemed to the young man, as he drew near and heard the fanfare and saw the waving flags, that such
a harvest festival must necessarily be destined to become a festival of betrothal. The representation of this became so
vivid to him that he quite began to consider how he should most fittingly reply when the toast of the betrothed couple
was proposed. But just as he rode in through the gate, the Unknown stood lying in wait and asked him in a whisper: are
you certain of your case?

% --- p. 4 ---
\setcounter{footnote}{0}%
\opage{4}How the omen struck him, the closer circumstances, how many painful moments he had to endure before he
reached certainty: on such particulars we shall not dwell. At long last, in one of the garden's dark corners, he was
overtaken by the young lady; there was no time for explanation, the crowd swarmed around them, there was a rustling in
the branches, there were people nearby: only this did she manage to utter, with a voice choked with tears and an
expression of the deepest pain: All is over! with these words she vanished like a vision. --- ``All over!'' Had an
abyss opened, had the comet lashed the earth with its tail, the young man could not have been more dismayed.

``All over.'' Like an outlaw who anxiously shuns the company of men, he hurried along a secret path through the
thicket, over a stone stair, out onto the open field. The voices and the glad laughter faded away, the tones of the
music sounded ever fainter and died softly away in the moonlit night. All was still; only a few distant barkings of
dogs were heard from the surrounding villages. He waded through loose earth, ever forward, as if he had a goal, until
he stopped at a large marl pit. He measured it with his eyes; the water, he knew, was three times as deep as the steep
edge was high, and the edge exceeded a man's height. The Unknown whispered: ``You must away, you must travel.'' He
looked into the mirroring deep; the comet stood down there, as if it were holding still and waiting for him. It is a
bird of paradise with golden feathers; with it I will travel. Upward I could not reach it, it sat too high for me, but
downward, a leap, a bold leap into the deep, and I ride on its back. But what is that black thing over there? The body
of a dead
% --- p. 5 ---
\setcounter{footnote}{0}%
\opage{5}dog! A travelling companion who in his haste has left his overcoat in the lurch: so I shall not have to make
the tour alone after all. ``You must first take leave'', whispered the Unknown; ``the journey is long, who knows when
you will come back''.

And he would take leave, not of any living mother's soul, but of the dark grove on the hill yonder. Up there was a
bench of turf; there slept in the grass among withered straws a sowing without harvest, a first-fruits of a heart's
spring, forsaken children of a disappointed hope, memories of a brief but indescribable bliss. So, he thought, first a
farewell in the grove, and then back to the marl pit and the comet and the dead dog! He hurried to the hill, there was
fever in his blood, his brain burned; he ran against a young but withered tree; an owl screeched, a startled hare darted
close past him. Everything spun round, he had a sensation as if fire were raining from the comet down on his head. With
tottering knees he worked his way up the footpath and caught sight of the bench. There sat a female figure. Is it
possible? she is sitting up there, she is waiting for me? No, it is impossible. Yes, it is she, and yet --- not she.
The seated one rose, lit by the moon's rays, and the figure became spirit-like, majestic. ``You are going away from
here, you would take leave; but you shall not travel with the dead dog. A friend shall accompany you, a friend who
knows the way, a noble, faithful companion, reliable as that one who followed Beatrice's lover through Hell, through
Purgatory and through Paradise into Heaven''. Then the strange vision vanished; the young man was struck as by
lightning and fell to the ground. What afterwards
% --- p. 6 ---
\setcounter{footnote}{0}%
\opage{6}became of him, when he rose again, does not concern us. We have touched on his story only for the sake of the
Unknown, that is, for the sake of imagination; for the being that awoke to life in him when he awoke, that was in truth
imagination.

In order to come to insight into the essence of imagination, we had best begin with reality, its conditions and laws.
Reality plainly as it stands and goes is not the object of imagination; occupying itself with the real it leaves to
the understanding and the five senses. But reality conceals something that interests imagination in a high degree,
something that it everywhere spies out and discovers; and what is that, then? If we ask imagination, it answers us that
it is mysterious beings, invisible beauties, faces that peep out from behind the curtains of reality. It is then, in
other words, not reality itself, but its possibilities.

But the possibilities, the possibilities that set imagination in motion, are of very different kinds. We can see it
from our story. The old women's imagination sought the possibilities in the comet, but went astray in the laws and
misjudged reality. They were astrological possibilities, and the golden age of the astrologers is past. No one believes
in the influence of the stars when the question concerns earthly life, its destinies and conditions. Tycho Brahe's days
are out of currency like obsolete coins. And why? Because the positions, motions and mutual interactions of the stars
conform to laws that can be found out by calculation. It is otherwise with the kind of possibilities that beguiled the
young man's imagination; they do not let themselves be calculated as one calculates a comet
% --- p. 7 ---
\setcounter{footnote}{0}%
\opage{7}orbit; they are psychological possibilities, possibilities of freedom.

The stars in the outward heaven follow such invariably determined orbits that the observer can state exactly in
advance the place they will reach at any fixed point of time. But the stars in the inward heaven, where each single
human being must preferably look for his lucky star, move in irregular orbits. These too are indeed subject to laws
and natural conditions, but the circumstances vary, and the perturbations are so manifold that no calculation is
sufficient: the understanding works itself weary, and imagination laughs at the embarrassment of prudence. Our young
man went boldly forward, he believed in his lucky star and was disappointed; he could just as well have been
disappointed in an entirely opposite way, through reserve and mistrust.

Among the Diapsalmata that introduce S. Kierkegaard's ``Either/Or'' the following occurs --- in what he calls an
ecstatic lecture, a lecture that our young friend perhaps has not had time to hear: ``Believe a girl, you will regret
it; do not believe her, you will also regret it; believe a girl or do not believe her, you will regret both; whether
you believe a girl or you do not believe her, you will regret both.'' It is a schema that can be filled in in the most
various ways. From ethics, for example: Be honest, you will regret it, be dishonest, you will also regret it; whether
you are honest or you are dishonest, you will regret both. From the intellectual sphere: Use your understanding, you
will regret it, do not use it, you will also regret it; use your understanding or do not use it, you will regret both;
whether you use your understanding or you do not
% --- p. 8 ---
\setcounter{footnote}{0}%
\opage{8}use it, you will regret both. Thus one can, of course, go on to infinity. Now what is all this supposed to
mean? The author gives a short but pithy explanation: ``Hang yourself, you will regret it; do not hang yourself, you
will also regret it; hang yourself or do not hang yourself, you will regret both; whether you hang yourself or you do
not hang yourself, you will regret both;'' and then he adds: ``This, gentlemen, is the sum and substance of all the
wisdom of life.''

The irony of despair, which thus plays its dialectical game, can be illuminated from many different sides; a main
side, which we wish to bring out in this connection, concerns the relation between reality and imagination. To let
imagination rule over reality will not do; it would break loose, turn everything upside down and rummage about
dreadfully in the given order of things. But to exclude imagination will not do either. When what we in a finite
sense call the real is to be our one and all, then the possibilities of freedom are misconstrued, the treasures of
inwardness that lie hidden in the mind remain unused, and the springs of life dry up. Let us look at the dry man of
understanding. How does such a one go about it? In order not to be beguiled by illusions he begins by parcelling out
the possibilities into nothing but imagined cases and calculating the probable consequences. Now that may be very good,
but how far does he get? He sets up two cases, three cases, perhaps ten, twenty different cases. Now he is to choose.
But when he then goes to choose the one, it occurs to him that he could just as well choose the opposite; if he would
hold on to the one, another at once presents itself. He stands swarmed about by cases,
% --- p. 9 ---
\setcounter{footnote}{0}%
\opage{9}like one who is swarmed about by flies; there is something teasing in it; it is the possibilities of
imagination that play tag among the cases, they make a fool of him, they mock him just as the elves mock Master
Grimmemann by filching now his commission and now his spectacles.

To exclude imagination in order to grasp reality will not do; for one loses the eye with which one was to discover
the possibilities in a deeper sense.

But the matter takes on a far more serious aspect still when the question concerns the destiny of man, his position
and significance in the great visible world. Behind the visible lies the invisible; behind the sensuous the
non-sensuous; to all that is outer there corresponds an inner, and when we look more closely, it is always the inner
that determines its outer. It is, after all, precisely the inner that we call the essence of things, their
disposition, their nature. Natural science, the investigating science of experience, examines the things in nature,
their properties, forces, laws; it is assured that every thing must act in accordance with its nature; but what nature
is, that it does not know. It makes no secret of this; it distinguishes sharply between what it knows and what it does
not know.

Natural science goes still further with this distinguishing of its own; when it is a matter of the single thing or of
whole groups or systems of facts of like kind, then it understands that the domain of insight and knowledge can
continually be extended. There is much that we do not yet know for the moment; but we can come to know it in time.
When, on the other hand, the question concerns the first original, the source of things, nature itself, then
% --- p. 10 ---
\setcounter{footnote}{0}%
\opage{10}it stops at an impassable limit; then the sober and level-headed science of experience admits that from its
standpoint it cannot possibly, by continued investigations, attain to a knowledge of what things are in their inmost
ground. ``We know the properties of things, but not the things; the phenomena, but not their essence; we do not know
what force is, nor what magnetism is, what electricity is, what life is and whence it springs, But this we know, that
the whole develops conditionally in necessary connection, and that everything must go according to unalterable
laws''.

Now all this would be exceedingly fine, enlightening and guiding, if we could remain standing here. But there is a
catch in the thing, and a sting in consciousness. We speak of an inner, invisible world that reveals itself in the
outer and sensuous one. From what we see, we infer to what we do not see. Now how does it stand with this inner,
invisible, non-sensuous world? Is all that is inner necessarily absorbed in its outer, or is there a surplus here? Is
the outer, in which the inner comes forth, confined to this world, this nature, or is there a higher world, a deeper
nature, of whose laws, life and forces this sensuous, visible world of ours gives only an imperfect reflection: can
there, in other words, be a supernatural world? The consciousness of culture, with ``modern science'', declares
outright: it is impossible! The supernatural is semblance and imagination. The heaven where the Almighty is enthroned,
the hell in which the damned are tormented, the angels of light, the demons of the abyss, all of it. Semblance and
imagination!

Here is a turning point, decisive for the course of development and the future of the generations. A fundamental
question is
% --- p. 11 ---
\setcounter{footnote}{0}%
\opage{11}pushed back, mistreated, bungled, and comes forth only in confused representations; it is the question of
the essence of imagination and its significance for the knowledge of the eternal truth.

The supernatural, it is said, is only an imagining. In this general proposition there is something that we must
concede, and something that we must combat for life and death. We must concede that all the representations we can
form of a world beyond, a supernatural realm, are formed in the form of imagination, that in so far they all belong to
imagination. But only an imagining: this little word only is the cardinal point of the dispute, the center of the
battle, this only decides victory and defeat, this apparently insignificant only is like the ``old farmstead'' of the
battle of Waterloo.

By this only, discord is set between reality and imagination. Here, then, what is thought of is exclusively human
imagination in its manifold, partly involuntary, partly voluntary forming, its fermenting and weaving, working and
shaping. It is then quite true that human imagination, with all its creative capacity, stands powerless over against
real nature; it can play with natural possibilities; it can sit at the loom of probabilities and weave the invisible
threads into patterns that have a deceptive, striking likeness to the real; but reality itself it cannot produce, it
cannot create out of nothing. But what if on closer investigation it turned out that not only the form of thinking but
also the form of imagination is an original and necessary form of truth; that there is not only a human but also a
divine imagination, not only a powerless but also an almighty
% --- p. 12 ---
\setcounter{footnote}{0}%
\opage{12}imagination, not only an imagination playing as in a fairy tale, but one willing in earnest, inspirited
throughout by the laws of wisdom, by plans and purposes? What if all that we call the real, the tangible, turned out,
with all its realness, to be not something absolutely self-subsistent, but an image stamped by creative, intuiting
thinking, an outcome of the power of divine imagination: that would change the matter. Could this be clarified,
grounded and made evident to the cultivated consciousness, would not the supernatural then be reconciled with the
natural, and the whole of existence present itself to us in a new light?

There is no talk here of a breach of the laws; nor will the mystery impenetrable to all human knowledge become
transparent: mystique and mysticism we must keep away. But a new depth of possibilities will open before us. In order
to look down into this depth without growing dizzy, we must stand firm on the foundation of the strict laws. It is a
new heaven and a new earth that we are to discover, not in the form of reality, but in that of possibility. Since
everything is to be seen under the veil of possibility in the mirror of imagination, it is necessary to know the
boundary line along which we are to move: to clarify the essence of imagination is not the same as to fantasize. We
must have the plumb line of reality with us, in order to test how deep the soundings are in the ocean of
possibilities. It might then turn out that there can be worlds of another and higher order than those that belong to
this natural universe.

We cannot follow Beatrice's lover on his lofty flight; we must go the level way of investigation and use our measuring
lines; we must distinguish between illusion and reality, between semblance and truth. Poetry has
% --- p. 13 ---
\setcounter{footnote}{0}%
\opage{13}its imagination, but religion has its own too. Imagination governs men just as well mediately as immediately.
We shall attempt a journey through the great realm where the Wondrous One dwells and reveals her power; this realm is
not far away; on the contrary, it is infinitely near to us. Many of the phenomena of imagination are known to us; but
the essence is an unknown: we shall strive to make closer acquaintance with this Unknown.

% --- p. 14 ---
\setcounter{footnote}{0}%
\lecture{Second Lecture.}\opage{14}

There is in our existence scarcely any subject of higher rank that has been able in such a degree to draw the
attention of the great majority to itself as the question of the immortality of the soul. It is a perpetual question
of dispute. The dispute is waged not only outwardly between those who decisively affirm and those who just as
decisively deny immortality; but the subject of dispute steals into each human being's inner self; it awakens discord
between sense and understanding, between head and heart, between faith and doubt. It is a persistent question of
anxious care; it fastens itself in consciousness, it penetrates through blood and nerves into conscience. How
sensuously dull or how thoughtless and frivolous must that human being not be whom this question of the soul and its
immortality never touched. It is more than an ordinary question of thought, % print: Tankespørgsmal (misprint)
it is personal. From this subject of dispute it is impossible to exclude imagination; it meddles in the matter whether
one notices it or not; indeed it ends by involuntarily coming to play a leading role. Let one arrange the text as one
will, it is still imagination that writes the notes and sets the music % print: Nodern (misprint)
% --- p. 15 ---
\setcounter{footnote}{0}%
\opage{15}to it. No theme has more variations. How much fantasizing there is over the immortality of the soul!

To be sure, the question of the immortality of the soul originally coincides with the question of eternal life, and in
so far is to be answered religiously in faith; but faith too has its imagination. Between faith and imagination there
is a peculiar relation. Faith will just as little as imagination subordinate itself to the outwardly sensuous; both are
agreed not to crawl before the laws of nature or bow the neck under blind necessity. Both feel themselves freeborn,
able to take up the struggle with the finite and rise above the barriers set for them by material reality. And yet
faith cannot without misgiving enter into league with imagination. It senses a danger. There is in imagination a drive
to go adventuring, a youthful wantonness, a desire to tease reality and play a game of probability with the great
possibilities. Into this game faith cannot enter, for its essence is earnestness; faith can, if need be, defy the
temporal world, for it knows its perishableness and has seen through its nothingness; but reality itself it cannot
defy; on the contrary! it is infinitely interested in the real, the higher real, for which it fights. Faith dare not
cast the burdens of reality from it in order to make itself light and soar high on the wings of imagination; for they
are wings of Icarus, wings of wax, which melt in the fire of the sun.

But, seen from another side, faith and imagination will make an alliance, namely where it is a matter of coming to
terms with the probable and recognizing the possible. Faith can go against probability; but it
% --- p. 16 ---
\setcounter{footnote}{0}%
\opage{16}cannot go against possibility. The improbable can be believed, but the impossible cannot be believed; so they
must both, faith and imagination, help each other to clear up possibility.

If what has been said here is applied to the doctrine of the immortality of the soul, then the first and most important
thing is surely this: whether such an immortality is possible. This question becomes ever more pressing the more we
come to know this real world, present nature. For antiquity the question was not nearly so difficult or involved as it
is for the present; for the present has a far more penetrating knowledge of the laws of nature and of the significance
of the material than antiquity had. There are in our time people who live in secret anxiety over the growing influence
of natural science; they have a fear of engaging with its investigations, they are afraid of losing their faith. It is
a painful condition and a sad testimony that all is probably not quite right with this faith, that there is in the
heart a worm of doubt that gnaws. If faith has truth, does it then need to live in a twilight? it must be able to bear
the light of day. There is no danger at all in letting natural science advance with its objections; it will surely
turn out that it is itself in error and oversteps its bounds when it fancies that it can blot out possibility, the
possibility in which faith is grounded. Should it seem for a moment to succeed, that is only by virtue of an
abortive thinking and a misguided imagination.

We shall cite some of the best-known objections to the immortality of the soul. It is impossible, it is said, that the
soul should outlive its body, for the life of the soul is dependent on organs. What we call our
% --- p. 17 ---
\setcounter{footnote}{0}%
\opage{17}soul is not a being by itself. The soul is indeed not matter; but it is nevertheless an activity of the
material, an activity of the nerves and above all an activity of the brain. Where does the soul go when it leaves the
body? Well, where does light go when the ether comes to rest; where does the electric current go when the battery is
disturbed; where does fire go when the chemical process ceases? It is with soul and body as with the tones and an
instrument. An instrument out of tune gives false tones, and when the strings snap, it is all over with the music. It
is a childish representation that the soul should leave the body or, as it is said, depart out of the body. A naked,
bodiless soul is an empty imagining. In the Orient they believe in the transmigration of souls; the souls of the
departed can dwell in plants and animals, in cats, dogs, elephants; indeed the believing Hindu must even be afraid that
he might meet the spirits of his forefathers in rats and mice: anxious fancies, uncanny imaginings. Among us, belief in
the transmigration of souls is indeed among the rarities, but all the more common is the representation that the soul
passes over into another world. Now where does that world lie? Imagination falls into embarrassment when it is examined
by the understanding. It invents a world beyond, but where is this beyond? Behind, outside, inside, below, above the
present world? These questions imagination cannot answer. Wherever it points, upward or downward, to north or south, to
east or west, it points everywhere out into space and does not get outside this natural universe. There is
fantasizing about a new body, a spiritual, supernatural body, hence a spiritual, supernatural
% --- p. 18 ---
\setcounter{footnote}{0}%
\opage{18}matter. But if there is anything natural, then surely it is matter. What, pray, is a supernatural matter? An
unnatural nature, an extended being without extension. For that reason, in recent times, people have accordingly begun
to arrange ``The Day after Death'' in a manner somewhat more in agreement with natural science, by making use of
hydrogen and the ether and letting the soul depart on a ray of light to other globes, the Pleiades or the North Star or
perhaps the Dog Star. It smacks of science; but the taste is insipid, and only slight reflection is needed to convince
oneself that the whole thing is a loose fantasticality.

When these objections are developed and carried through with keen understanding and knowledge of the laws of nature,
they make a strong impression. The further deliberation goes into particulars, the more absurd the whole becomes. The
true and the false are mixed in so peculiar a way that the matter can easily take on the appearance that faith in
immortality is only an imagining, a fancy. Now whence does that come? It comes from the fact that the foundation on
which the whole view rests is false.

It is, in the first place, a false, a fundamentally false representation that the soul is not a being distinct from
the body. In order to prove that feeling, consciousness, understanding, imagination are nothing but activity of the
nerves and the brain, one proceeds from experience; but experience comes into conflict with itself; for there is not
only an outward but also an inward experience. I experience that I have a body, and that this body consists of material
parts; but I experience also that there is something in my inner self that I must call myself; I experience that my
consciousness, my feeling, my understanding and my imagination,
% --- p. 19 ---
\setcounter{footnote}{0}%
\opage{19}however much these activities are conditioned by and dependent on material organs, on nerves and brain, still
cannot possibly be material activities. The magnet can attract, an electric battery can develop electricity, Such
things belong to the nature of the lifeless elements. But matter, however peculiarly it may be distributed and
arranged, cannot feel; thinking is done by means of the brain, but the brain, this strangely furrowed mass with its
convolutions and labyrinths, the brain as merely organized matter cannot think. It is our soul that feels, it is our
soul that is conscious of itself, it is our soul's being that thinks. Now can it do any good to come with the
explanation that the soul is only an imagining. Whence, then, comes imagination? We should like to hear the physicist
who could explain to us with any plausibility how the atoms must arrange themselves in order that matter may begin to
fantasize. Let one try to live for a day, or even just an hour, with the view of the materialists or the modern
positivists, and one will gradually approach madness. If the soul is to be brain and the brain soul, then the whole
human organism must be regarded from the point of view of a peculiarly formed heap of organs; one cannot then speak of
my body or your body, of my brain or your brain, for expressions such as mine and yours presuppose that you and I are
particular beings. When two human beings talk with each other, it is thus not two souls but two brains that
communicate their representations to each other and confide their thoughts to each other. To an ironic humor it may
indeed be quite entertaining to regard people in the squares and streets as self-moving heaps of atoms, organic
automata, to see them nod to each other and take off their hats to each other,
% --- p. 20 ---
\setcounter{footnote}{0}%
\opage{20}while each of these brains fancies that it has a self, an I, although all that about I myself and you
yourself is only imagination. In one of Hoffmann's novels a young girl is depicted who by her beauty and her talents
makes an extraordinary success in elegant society; she dances gracefully, she sings and plays with virtuosity, with
taste and expression. The young gentlemen flock around her, she captivates them in a magic circle; there is adoration
and courtliness, there is heart-throbbing and rivalry; there are flashing glances that even threaten duels. But one of
the adorers keeps quietly back, he waits. He waits for the moment when he can meet the beauty alone and throw himself
at her feet. The moment comes, but what must he learn? The adored one is no real human being; she is a doll, an
extraordinarily finely wrought work of art! The eyes of polished coal, the dark brows attached, she is an automaton,
she has no soul, she runs on screws. What if we were all now automata produced by nature, who made ourselves and others
fancy that we had a living soul, while the truth nevertheless was that we all ran on screws. Madness!

It must stand firm for sound sense and inner experience that we have a soul, and that the soul active in our body is a
being distinct from matter; and then it is not, as the materialists consider it, childish to speak of the soul leaving
its body at death. But where does it go? Here imagination comes to have a word to say, and that in earnest. What is
space, what is matter, what is the bodily world? Without the help of imagination it will
% --- p. 21 ---
\setcounter{footnote}{0}%
\opage{21}be impossible to give any intelligible, let alone any reasonable, answer to these questions.

What is space? a form for the extended, in itself an empty form, but nevertheless a form of intuition, a form for all
that is intuitable. This intuitable can now be partly something outward and material, partly something inward and
spiritual; but space always conforms to the content with which it is filled. He who dreams, dreams in space; but it is
not outward space, it is an inward one, and yet this inward thing takes the form of an outward one. The traveller far
from home who is longing, for example, can lie in a strange place in the distant, foreign land and dream of his home.
The dream can be so vivid that he sees them all before him in the familiar rooms and the well-known surroundings. There
is thus another space in his dream than the one in which he bodily is, and in which he will find himself when he wakes.
Here, then, are two kinds of space, the dreamed and the real. Now why do we not call the dreamed space real? Because it
is not filled with real objects, but only with dream-images. As the empty form, the dreamed space is just as much space
as the real one: the difference rests solely on that with which the space is filled. A painting on the wall, a
landscape for example, shows us the same thing in another way. The space in the room and the space in the picture give
two kinds of space. The picture, whether it represent the Roman Campagna or the Jutland heath, has its own space.
Foreground, middle ground, background separate; the nearer draws near, the farther draws away: we look into an
indeterminable depth, and yet the wall is not broken through. Now what kind of space is that? An
% --- p. 22 ---
\setcounter{footnote}{0}%
\opage{22}illusory space, and why? Because the objects with which it is filled are illusory.

Imagination, then, is not in contradiction with itself when a distinction is drawn between different kinds of space.
The whole difficulty about here and yonder falls away as far as space is concerned. We must then ask exclusively about
that which fills space. We thereby come to gain a deeper insight into the essence of imagination.

Our human imagination can fill an inward space with images, but the space remains illusory, a fancied space, because
our images are fancies that lack reality; it is the powerlessness of our imagination that it cannot give its images
true realness. What, then, is required to produce a real space? An almighty imagination that can fill its space with
real objects, a creative imagination that is able to translate its rich content, its inner, into a manifold and mighty
outer, that is, into an outward, sensuously existing world.

It is the relation between the inner and the outer on which everything here depends. We speak of the Creator's
almighty thoughts and of his almighty will; but let us not forget his almighty vision. If in the Creator's eternal
self-power there were nothing but thoughts and will, then this outward, bodily world would never have come into being.
There are indeed thoughts, divine thoughts, in this world, but here there is also matter and corporeality; this world
of matter and corporeality is more than a world of thoughts. Thoughts are invisible, but here is a world of tangible
things. In the laws of this world a law-giving will is indeed recognizable; but there is more in the world than laws,
there are
% --- p. 23 ---
\setcounter{footnote}{0}%
\opage{23}substances and masses. The laws are invisible like the will from which they proceed; but the order of things
in which the laws are fulfilled and the will accomplished is sensuous, visible, tangible. Now whence stems this outward,
sensuous, tangible thing? From imagination, the almighty imagination. It is, after all, precisely the essence of
imagination to mediate between the inner and the outer, the invisible and the visible. Imagination is akin to thought
and belongs to the inner; but it is also akin to the sensuous, especially to the sense of sight: its productions are
images, and when there is power in these images, then the images are real objects. Power and reality, power and
self-subsistence, power and realness correspond to each other in the least as in the greatest. There is a
self-subsistence in things because there is a power in things; there is a self-subsistence in each atom because there
is something of a power in the atom: the more power, the greater self-subsistence; the less power, the less
self-subsistence. The whole of nature is a realm of power. Therefore nature has a closedness in itself, it develops
out of its inner, its dispositions; but all determinations are given in advance, it cannot give itself a single new
disposition. This power in nature stems from omnipotence; it is a making-visible in divine imagination, a manifestation
of the infinite, creating and sustaining self-power.

How many worlds can there be, then? We answer: as many as can be contained in the divine imagination, and that is ---
many!

The Savior's words: ``In my Father's house are many mansions'' indicate that his gaze, even if it
% --- p. 24 ---
\setcounter{footnote}{0}%
\opage{24}perhaps rested on the oriental starry heaven, on the Pleiades or Orion's belt, nevertheless penetrated this
outer and looked into another heaven, a world of higher order. There can thus be, besides this visible universe of
ours with its infinity, its space, its matter, its corporeality and its countless globes, innumerable higher worlds,
each with its own infinity, its space and its corporeality. Our starry heaven is high and deep: but deeper still is
the almighty imagination with its inexhaustible possibilities. From this point of view we shall consider the
immortality of the soul.

Now what does it mean that the soul leaves the body? We express it most clearly by saying that it exchanges the one
world for the other. In order to see what lies in this, let us cast a glance at sleep and death. Death, it is said, is
the brother of sleep; the genius of sleep with the upturned torch corresponds to that of death with the downturned one.
This brotherhood does not hold in the physiological sense; the physicist and the physician would have much to object to
it; but it holds in the psychological sense; it holds when the question concerns the inner, the essence and ground of
life. In sleep the outward senses close, consciousness diminishes, and sleep can be so deep that even the dream
vanishes. In death all sensing ceases, the light of consciousness is extinguished; but the soul does not become
nothing; it sinks in its inner down into its deep ground; this ground is the eternal being, the divine imagination, in
whose bosom it rests. Therefore we name the dead with the significant word, those fallen asleep. Upon the sleep of death
follows an awakening, a new formation of a body in a new world. The new world is of another and higher order than this
present one;
% --- p. 25 ---
\setcounter{footnote}{0}%
\opage{25}but both the world that the soul leaves and the one into which it enters by the formation of its new body,
both worlds, however different in kind they may be, are manifestations of the one and same almighty imagination. With
the awakening follows the remembrance of what is past; but the transition is no journey, it is a transformation. The
transition is inward, and yet real; we call our dead the departed.

So far we can go without running up against any whatsoever of the known laws of nature; but if we venture further
forward, if we attempt to depict how things look in the new world with the new conditions, then we begin to fantasize.

There is in certain people an unbridled, curious craving to make intimate acquaintance with the world of spirits,
especially when it can happen in a properly sensuous and tangible way; they conjure up spirits, they question the
dead. It went far in antiquity, think of ``Emperor and Galilean''; but it goes much further still in the present.
There is in this respect no undertaking so foolhardy that an excited imagination will not dare to engage in it. To get
the spirits to show themselves bodily, especially in the dark, is now as nothing; but to interrogate them quite
privately, to photograph them, to take them into the household and get them to make sausages, that is --- ``a
science''. It is revolting when such a fantasticality of the grossest crudity of spirit wishes to put itself in
connection with religion and help to get Christianity scientifically grounded.

There is a form of imagination that stands near to religion; it is of nobler, higher origin, it has its
% --- p. 26 ---
\setcounter{footnote}{0}%
\opage{26}source not in gross sensuousness, but in a freely appropriated, clear intuition of the crises of life; it is
a poet's imagination that unfolds its content in so-called apocalyptic poetry. Here belongs, as a member of world
literature, Dante's divina Comedia, in our own literature Heiberg's ``A Soul after Death''; the concluding part of
``Adam Homo'' too should be mentioned. The images that occur in such poems make no claim to count as reality, they are
symbols, They can be of gripping effect, and acquire ethical significance by making vivid states of the soul that
remind us in a penetrating way of the personal laws of development and bring to life our representation of the
connection between life here and life yonder. Think of Adam Homo as he sits on the sinners' bench and looks with
horror at the great scale. Alma approaches, a star of hope in the dark night; she comforts the disheartened one, leads
him by the hand and inspires him with courage to go through the ordeal by fire. It is a struggle between the new and
the old, a process of transformation that is made vivid in a situation. The formerly great one becomes so small, and
the formerly small one so great. Adam Homo, in the world the great man, ``Adam Homo with the gentle spirit, Baron,
Privy Councillor, Knight of the White Ribbon'', how small he has become; and Alma, the quiet, deep, but overlooked
% „Adam Homo med den blide Aand, Baron, Geheimeraad, Ridder af det hvide Baand“
and forsaken soul, how great she is! In order that the new glory can be revealed, there is an earthly substance that
must first be consumed, an ``armor'' that must melt, an old glory that must be ``burned to coal''. And it takes place
before our eyes, it becomes present to us, it sounds forth from the chorus of Purgatory: ``The heart's melted armor
flows away, With
% --- p. 27 ---
\setcounter{footnote}{0}%
\opage{27}our finery the flames make sport: Our wreaths, exploits and virtues go up as a nothing in smoke''. It is
% „Hjertets smeltede Pandser henflyder, Med vor Pynt drive Flammerne Spøg: Vore Krandse, Bedrifter og Dyder gaae op som et Intet i Røg“
imagination, but this imagination is under the discipline of the idea, it is no idle, empty fantasizing.

And yet there is a difference, an essential difference, between apocalyptic poetry and the word of revelation. The word
of revelation makes vivid in broad strokes what awaits us in a world to come. These depictions too belong to the
domain of imagination; but here a quite different standard must be applied: here what is aimed at, in a higher sense,
is something real. Here there is no basis for further-reaching fantasizing, but a summons to earnest self-examination.
In what follows we shall pass on to consider religious imagination.

% --- p. 28 ---
\setcounter{footnote}{0}%
\lecture{Third Lecture.}\opage{28}

Already a good many years ago attention was drawn in our literature to the fact that religion was suffering mistreatment at the hands of science. It was then emphasized that a distinction had to be made between the natural and the supernatural, that manifestations from a realm beyond could not be measured by the yardstick of everyday experiences or pressed in under the known laws of nature. It was emphasized that, although nothing can be true in religion whose untruth can be proved scientifically, it was nevertheless improper to set science up as judge over religion; that was to idolize our piecemeal knowledge and raise it up onto the throne of omniscience. Religion had to be allowed to develop and ground its truths religiously; then science could see to it that nothing was contained in the development that conflicted with notorious facts, and in any case exercise its indisputable right to free and many-sided investigation. It was shown in detailed exposition that, in consequence of the limitation of human knowledge, there must necessarily be two circles of knowledge, independent of each other and yet mutually illuminating each other, one for the supernatural and one for the natural, one for religion and one for
% --- p. 29 ---
\setcounter{footnote}{0}%
\opage{29}science. In short, an unambiguous separation was made between faith and knowledge.

This separation gave offense, not only to those who held themselves assured in advance that the days of faith were numbered, since the supernatural, regarded in the light of knowledge, is only semblance and imagination; but, what is more remarkable, also to those who felt themselves called, on behalf of the doctrine of the Church, to plead the cause of religion. They declared outright that the supernatural not only could be taken up into science as a ground, but that there existed a science supernatural in its origin, a science of the Holy Spirit, which enabled its devotees to raise themselves above ``the rudiments of the world'' and chastise ``the heathenishly puffed-up natural sciences.'' And what way did one then go? Not the way of logic, for fear of infecting oneself with ``the Kenological''; one chose a more graphic procedure, namely to have such words as will, Creator's will, free inwardness, the good will, printed in bold type. Bold type is, after all, a reality too; and with such realities one then prepared oneself to chastise the conceited real sciences. As was said, the offense at religion's having its own circle of knowledge was so great on both sides,
that the defenders of faith and the deniers of faith, men of the Church and men of free thought in the name of science, gave each other their hands, and the mutual understanding seemed about to become cordial. The cordiality faded, however, and the whole passed little by little into a quiet ferment. The honored general public has since that time been anxious now for science and now for religion, now for ``faith'' and now for ``free thought''; but
% --- p. 30 ---
\setcounter{footnote}{0}%
\opage{30}in the last few days the prospects have grown brighter; the proper parties, who are thoroughly at home in the matter, have at last taken upon themselves the arduous task of enlightening the great public. The anxious general public will now be able to inform itself of how things really stand with the disputed account between religion and science. That a decision is approaching is evident; it can clearly be seen from the fact that the dispute is being carried on in the newspapers. When such a dispute about such a matter passes into the newspapers, one knows that now it is in earnest; it is just as unfailing a sign of earnest that a matter of this nature gets into the newspaper as it is an unfailing sign that a melody is generally liked when it gets into the barrel organ.

We cannot here take up the whole matter with the brevity, the pith and the clarity with which it ought to be taken up in the newspaper; such things we leave to those who have the gifts, the preliminary studies and the authority for it. But one contribution we can nevertheless give, namely towards illuminating the side that especially concerns imagination.

Revealed religion is grounded on miracles; but --- for this ``modern science'' vouches --- the miracle is only semblance and imagination. No miracles happen and no miracles have ever happened; the miracle is a self-contradictory and therefore meaningless representation. The miracle is a breach of the laws of nature, but the laws of nature are unchangeable; that is a rigmarole which every schoolboy nowadays has at his fingertips. --- But what, now, is really meant by the laws being broken? What is meant is that the connection of things is broken, that the divine wisdom contradicts itself. This proposition holds on one condition, namely that
% --- p. 31 ---
\setcounter{footnote}{0}%
\opage{31}the whole plan of wisdom is exhausted, completely exhausted, in this natural order of things. Can science really prove this --- as yet it has not proved it; its alleged proofs are only assertions, and these assertions are in poor agreement with its sober confessions that it does not know the essence but only the properties, not the force itself but only its manifestations, etc. --- but should science one day attain so deep-reaching an insight into the essence and secret connection of things that it can demonstrate with clear grounds that the supernatural world is only semblance and imagination, then it is truly all over with revelation and faith. But not before; let us not be hasty!

The deniers of revelation are clever enough to keep the dispute fixed in the domain of the laws of nature and to act as though the question of the supernatural had long since been settled: in this way they already have the victory in hand. What good does it do that the defenders of faith appeal to the almighty Creator, when the Creator's plan, with laws and connection, matter and forces, is exhausted, completely exhausted, in this natural universe? If there is nothing behind it, nothing at all, no free inwardness, no depth of riches of possibilities from which creative freedom can draw, what use is it then to call upon omnipotence? The eternal law is the lawgiver's master, necessity reigns: the pulse of freedom stands still. The human will does have a kind of free inwardness, a self-determining; but if the bosom of the Godhead is immediately one with the ground of nature, then everything breaks forth, yes, life itself, out of conditioned necessity, as when a spark falls into gunpowder. We live in the age of
% --- p. 32 ---
\setcounter{footnote}{0}%
\opage{32}constitutions, and science has had its eyes opened to the constitution of the great whole; it will hear nothing of ``an unaccountable God'' with unsearchable counsels; it has gone through the old fundamental law of existence anew and assured itself that the Almighty will surely refrain from making a coup d'état.

Let that be imagination; the question nevertheless remains whether the defenders of faith appointed by the Church are for their part really able to introduce the supernatural into science and get the knowledge of it sufficiently grounded. To put on the air of authority and use the language of power does not help; to want to force an absolute monarchy upon the free republic of the sciences, and to have ``the good will'' printed up in ever so bold type in order to chastise the naturalistic rudiments, leads to humiliating defeat. Here is a secret of imagination which we must see to getting cleared up.

Revealed religion has visions and images, wondrous forms of creative imagination; mind you, not merely of human, no, originally of divine imagination. The long, thousand-year dispute about religion and science bears witness to this. The believing Fathers appealed to the testimony of the Spirit, and rightly. But they understood the testimony one-sidedly, in that they did not sufficiently heed religious imagination, and therefore overvalued the outward, the historically real. That served for a time. But then came a turning point. As the enlightenment of the understanding increased, and criticism began to come into its strength, difficulties arose, the one more intractable than the other. The Protestant-orthodox teachers, who in contrast to the Catholics believed they possessed a reliable, juridically
% --- p. 33 ---
\setcounter{footnote}{0}%
\opage{33}registered document of revelation in Holy Scripture, proceeded from the presupposition of the absolute infallibility of Scripture. To glorify that infallibility they assumed that every word, every letter and tittle, in the Old as in the New Testament, was breathed in by God; this assumption prepared painful disappointments for them. Infallibility required that the vowel points in the Hebrew text too must originally have been breathed in; but criticism showed that these signs were a later addition by Jewish scholars, and the spokesmen of orthodoxy had to give way. Infallibility required that the Greek of the New Testament must be classical, for the Spirit that had in supernatural fashion distributed its gifts of tongues at the feast of Pentecost must surely be just as thoroughly at home in the Greek language as either Plato or Aristotle; but when philology put on its spectacles and came to look, this proved to be a mistake, and orthodoxy had to give way. So then, despite human imperfections in the linguistic respect, the whole biblical history from first to last was to be infallible in each and every part. But now the investigations took a new upswing. Learned researchers proved that many of the books were spurious; that the text was doubtful; that there were discrepancies, not only between the biblical accounts and the testimonies of profane history, but, what was still worse, deviations and evident contradictions among the sacred authors themselves. At last the objections rose to so overwhelming a flood that the old orthodoxy stood speechless, at a loss, and saw all its dams broken.

Then, as far as the doctrine of the Church was concerned, a more liberal and unprejudiced period set in.
It
% --- p. 34 ---
\setcounter{footnote}{0}%
\opage{34}was of no use, after all, to be stubborn and defend the indefensible; one had to bow to the proofs of science, one had to be reasonable. But gradually, as one began oneself to become reasonable, one also began to see that there was much in Holy Scripture that conflicted with reason; the devil, e.g., was not reasonable at all; the demons and the possessions and all the supernatural with which the written revelation was encumbered, were unreasonable. In the end one discovered that all miracles --- excepting that of creation --- were unreasonable.

Only the natural was real; the supernatural was unreal. One then had to see about getting rid of the supernatural; this could best be done by a philological-critical interpretation. Yet here one struck a stumbling-stone that was not easy to clear out of the way; for it was the cornerstone of the whole building, it was Christ himself. It was found regrettable that a man otherwise so reasonable as Christ had let himself in for performing miracles. The freethinkers openly declared him an impostor; but the teachers of the Church, who could well see that he was a paragon of virtue, excused him as best they could and gave assurance that he had done it with good intentions, merely to accommodate himself to the representations of the common people. Others shifted the blame from the Master onto the disciples; it was the disciples who had misunderstood him and transformed his natural healings of sickness into miracles. And what, then, became of the gospel history? A dry skeleton in distorting wrappings. The pure Virgin became an ordinary carpenter's wife, yes, a woman of doubtful reputation. The Risen One was a man only apparently dead, who, ailing from his wounds, crept
% --- p. 35 ---
\setcounter{footnote}{0}%
\opage{35}about in the twilight and got a splendid departure only insofar as he once, in a strong evening light, went behind a high mountain and thereby made the disciples believe that he had gone up to heaven. Poor it was, this revelation; but it was scientific. Then all at once a new way out appeared; it was a new vision: aesthetics and philosophy had caught sight of imagination, that is to say the human.

Whence come the heathen gods? They are not historical persons, not empty allegories of nature; nor are they arbitrary fictions which through priestly deception have found entry among the common people and founded their superstition. The heathen gods are creatures of an imagination filled with ideas but unconsciously poetizing. These productions in legend and poem are myths; religion is mythology. Now the hour of liberation struck. If the heathen religions are mythologies, why should the so-called revealed religion, in Judaism and Christianity, not also be mythology? Myth has precisely this peculiarity about it, that it springs forth in imagination with the irresistibility of a spiritual force of nature. Creative imagination borrows its elements from natural events and history; from this it is then explicable that the mythical can be mixed in the most naive way with the outwardly real, the supernatural with the historical. Now, then, one began to find myths, first in the Old and afterwards in the New Testament, nature myths, mythical legends, historical myths, philosophical myths, etc. And what, then, is the result? The newest scientific result is that all the supernatural in the New as well as in the Old Testament is mythology. The spirit, to be sure, is
% --- p. 36 ---
\setcounter{footnote}{0}%
\opage{36}gone, but one still has the idea, after all. On behalf of the idea the official doctrine of religion is for the present to apply criticism in order to find out what and how much in the Gospels is myth and what is real history.

The life of Jesus falls into three sections: 1) from his conception and birth to his baptism in the river Jordan and his temptation in the wilderness; 2) from his appearance as a teacher to his death and burial; 3) the third and last section comprises his resurrection and ascension. According to the scientific view now in force, only the second section, on his public life, contains a real history, and from the reality of this history his miracles must then be separated out, for the miracles are, after all, only myths. The first and the last section are pure and simple mythology. The annunciation a myth; the miracle of the birth a myth; the resurrection, with the ensuing manifestations, and the ascension, a myth. No one who wishes to stand on the ground of science is able here to maintain the historically real; it is not one thing or two things that make difficulties here; it is, if you like, a thousand, yes, ten thousand things, for it is the concurring testimony of all present sciences about the laws of nature and their unchangeability that the defender of the historical here has against him. It is just as impossible to get any one whatever of the many sciences convinced that these are not myths but real facts as it is impossible to convince astronomy that the heavens are a fixed vault and that the earth stands still at the center of the world. The individual may well assert it; but science pays no heed to his assertion. The defenders of faith can appeal to the fact that Christianity is a historical fact; science answers that
% --- p. 37 ---
\setcounter{footnote}{0}%
\opage{37}Mohammedanism, Brahmanism, Buddhism are also historical facts, without its following from this that their myths are divine revelations.

So, then, the old, traditional but spirit-forsaken doctrine of religion still stands on the ground of science. But how does it stand? Its position is not enviable; it looks with folded hands upon the great images of revelation and is seized with melancholy. The angel of the annunciation, the vision of the shepherds in the radiant night, how simply sublime. Would that it were real! But alas, it is not real, it cannot be real. Science says No! and the laws of nature threaten: so in God's name it will have to be myths.

If the matter were not so serious that the situation must in truth become tragic, the whole despair would really be comic; for it is self-inflicted: the whole thing rests on a fancy. A poet's imagination that wishes to acquaint itself with the newest scientific exploits of the official doctrine of religion will find an inexhaustible store of subjects for comic treatment. Now we are told that the Gospel of John is the gospel of the true, deep heart, and that the three others must conform to John. Then with the next post comes the news that the Gospel of John is a finely wrought, conscious work of art. Therefore not historical! Which of the evangelists is now to play the leading role? It can hardly be Luke now: therefore either Matthew or Mark. It is Matthew; no, it is Mark! Then one lights the lantern of criticism and searches for a primitive gospel.

Just as great is the perplexity over the historical in the miracles. One wants to distinguish critically between reality and legend, and to accept as real only that
% --- p. 38 ---
\setcounter{footnote}{0}%
\opage{38}which has the concurring testimony of all the evangelists. The walking on the sea, then, need not be real --- in the lecture chair, that is, for in the pulpit it is always real; but the miracle of the feeding must necessarily be real, for it has all the testimonies. But no, it cannot possibly be real: the loaves are too few and the fishes too small. ``To see the bread swell and multiply in Christ's hand is an uncanny sight: what will science say?'' --- Would it not then be best to exchange the scientific foundation for a religious one, and to choose one's standpoint outside science? ``No! no! that will by no means do; the day the doctrine of religion ceases to be a science, an objective science firmly grounded on objective facts, on that day religion must perish, and the sad tidings will stand in the newspaper. Nor is it necessary; the danger is not so great; it is only a transition. Science has not yet said its last word. One can be sure that when science once gets its last word said, we shall surely get the historical reality of the miracles objectively, scientifically grounded.'' But when does science get its last word said? It will be a long-drawn process, longer still than the one known from Bleakhouse: Jarndyce and Jarndyce. The human understanding is offended, and imagination bewildered; imagination, bewildered by offense, stands smiling among the documents and says with the old mad maiden lady: Pray, accept my blessing!

The doctrine of religion need not go out of the circle of the sciences; but it does in truth need to place itself on a non-scientific, i.e., on the supra-scientific standpoint proper to faith and
% --- p. 39 ---
\setcounter{footnote}{0}%
\opage{39}revelation. Then let the light of science fall clarifyingly over the whole.

If we place ourselves on religion's own original foundation, it will at once become evident how it comes about that myth and revelation can have so striking a resemblance to each other that scientific criticism, which has reduced the supernatural to semblance and imagination, is not able, despite all its acuteness, to tell them apart. They both have a common form, the form of vision, of imagination. The almighty God is also the God of vision and of visions. But there is a presupposition on which faith stands firm, and this presupposition no science can shake: it is the omnipotence of love. Science can neither deny nor affirm God's love; for it knows nothing at all about it. In the omnipotence of love lies the possibility of revelation; this possibility science will surely have to let stand, and then that is enough. In the possibility of revelation faith has a shield against the arrows of criticism; but in the actualization of this possibility the power of imagination is also revealed.

Love is a thought and more than a thought; it is a will and more than a will; it has a heart, and the heart desires a vision, and in the vision a figure, and vision and figure are seen in imagination. There is in myth something of a divine imagination; but the divine is darkened. Entangled in sensuousness, ensnared by selfish desire, human imagination left to itself has seen only glimpses of the truth, not the truth itself. The light has shone in the darkness, but the darkness has not comprehended it. Hence the veiled, the unclear and the ambiguous;
hence the fantastically
% --- p. 40 ---
\setcounter{footnote}{0}%
\opage{40}illusory with which myth is always afflicted. In revelation, on the other hand, divine imagination has power over the human, clarifies it in its light and fills it with its truth; therefore the phenomena of revelation are so pure, unambiguous, transparent to the spirit. Spirit and nature are inwardly united; the supernatural too has its nature. The concept of nature therefore has two sides, a scientific and a religious. The scientific concept of nature is formed, and must form itself for us, under the impression and observation of natural facts; but through the visible the invisible is glimpsed, through the known the unknown, through the laws the lawgiver, through necessity freedom. Spirit in its unchangeableness, with its thoughts of eternity, unfolds the content of nature in the connection of things, in the form of the necessity of imagination. Further than this science does not come.

Where science lets go, religion begins, and religion is the original; it goes before all science in time; it passes through revelation and myths. Here, then, it is a matter of making a separation between fancy and truth, of finding a touchstone that can divide revelation from myth. This touchstone we have in faith in the omnipotence of love. That the Creator of nature is also the Father of spirits, that the God of necessity and of the laws is also the God of freedom and of love, that the supernatural in free inwardness is the invisible foundation of all that is outwardly natural, this can be known only in faith. This faith, once planted in the ground of the life of the race and expressed in a definite, unambiguous confession of faith, science will never get rooted out of the hearts of men. The rigmarole about
% --- p. 41 ---
\setcounter{footnote}{0}%
\opage{41}the unchangeable laws becomes meaningless when it is recognized that there is not only a visible connection between the visible things, but also an inner, if invisible, connection between the natural and the supernatural; that, in other words, worlds of lower and higher order are united in the essence of omnipotence, in the divine imagination of freedom and of love.

The difference between revelation and myth can be seen, in the light of faith, by the thought, the will and the imagination. By the thought! There are thoughts in the myths, often profound thoughts, but the visions rise and fall, the crossing thoughts are mingled, and the representations confused. In revelation the thought is clear, the thought of the law, the thought of life, the thought of sin and death and the thought of restoration and salvation. By the will! Through the myths there runs a longing, a yearning and an urge to see a present God. But the sensuous and the spiritual, natural drive and selfish passions mingle with the deeper desire, and the holy will does not acknowledge it. So the wonders of myth become fancies, fantastic signs: here no miracles happen. In the word of revelation the will is present, the almighty will: thought and will correspond to each other as word and sign. The God of the spirit, the God of the law and of love, is present in his word and speaks in the sign the language of omnipotence: Hear the word, see here my sign! that is the miracle. But however outwardly real the miracle may be, it must still be seen with the eye of faith: thus divine imagination sets the human in motion. What a difference between the often brilliant, colorful, but, for all its grandeur and plasticity, nevertheless confused mythic imagination, and the clear, irradiated by the light of blessedness, and for that very reason so pure,
% --- p. 42 ---
\setcounter{footnote}{0}%
\opage{42}chaste imagination of revelation. If we cast a glance at the gospel of the annunciation, and let the divine thought and will and imagination be united in the omnipotence of love, then we can understand the revelation. Here is a thought that is to be introduced into human life, a thought of love, and the thought is this: since all, even the noblest natures in the race, bear in themselves from birth the seed of sin and death, so now One is to be born who by nature itself is sinless and pure, a man of the spirit, a child of eternity, a holy seed, which grows up without sin and saves from sin and death. The unity of this thought with the almighty will is the cause of the miracle. The miracle itself is a secret; but the secret is to be revealed, it is to become known to all; then comes the all-revealing imagination and makes the invisible visible. The greeting of the angels, Mary's words, the solemn declaration! in this image of life there is nothing mythical, nothing ambiguous; it says precisely what it is to say, neither more nor less. It is not human fancy; it is, in the omnipotence of love, divine imagination; every line and every tone of color, every single feature and every gleam of light upon it, all is of spirit; what is thus revealed is real in the very highest sense, but it is spiritual reality, reality for spirit. --- As the beginning is, so too is the continuation and the final close. But the lights alternate; the natural and the supernatural, the real outward and the higher inward, interlock wonderfully; only the imagination of revelation, in unity with faith, is able to make such an interplay clear to us: we shall see more closely how.

% --- p. 43 ---
\setcounter{footnote}{0}%
\lecture{Fourth Lecture.}\opage{43}

How easily a dispute about divine things can degenerate into an empty dispute about words, when due care is not taken that word and matter mutually illuminate each other, is well enough known. If there is a dispute, e.g., about faith and knowledge, one can very easily give it the semblance that all knowledge is grounded on faith. The painter, who is to use the colors, must, after all, believe, believe in his sight; the mathematician with his a = a must believe, believe in his understanding; the astronomer must believe in his own and others' observations, and so forth. But then from the other side one can just as easily give it the semblance that all faith is grounded on knowledge. We must, after all, know what faith is, know what it is in which we believe, and know why we believe in it; and when we know all this thoroughly, the whole is, after all, built on knowledge. This sham fencing stops when we honestly go into the matter. There is something that is not to be believed but known, and there is something that is not to be known but believed. If one says that the sun is shining, the man who sees answers: I know it; only a blind man could say: I believe it. If one says that two and two are four, the man of understanding answers: I know it; only an idiot could say: I believe it. If the talk is of a sum of money, and
% --- p. 44 ---
\setcounter{footnote}{0}%
\opage{44}the question is asked: is the account correct? then it does not help to answer: I believe it, for it is not to be believed, it is to be known. In such connections faith means only a provisional assumption, a half-certainty joined with doubt. Here knowledge is sought; only with knowledge comes certainty, and with knowledge and certainty faith disappears. He who in such cases has attained full certainty says: I do not believe it, for I know it. It is otherwise when the talk is of the immortality of the soul, of God's love, of the supernatural world. A human being can, as regards all this, be so sure of his cause that he dares to say without circumlocution: I know it; I know for certain that God is love; I know for certain that there is an eternal life, a resurrection, etc. Knowledge here means an unshakeable certainty; but this unshakeable certainty does not make faith superfluous, on the contrary! it is one with faith itself: knowledge is here an expression for full assurance in faith. From this point of view too faith can be set against doubt, but doubt then does not have its source in imperfect knowledge, but in unbelief. What must be believed in the very most proper sense is precisely that which by its nature cannot be known; --- only it must not be overlooked, but must be most emphatically impressed, that faith in its sphere of knowledge can give an equally reliable certainty, an equally perfect assurance, as knowledge in its own.

In the sphere of faith the opposition between faith and knowledge has vanished, for here there is no question at all of knowledge in the scientific sense; here oppositions of an entirely different kind enter: the opposition between faith and superstition, between unbelief and superstition.

To come to clarity about how matters stand with the oppositions in the domain of faith,
% --- p. 45 ---
\setcounter{footnote}{0}%
\opage{45}there is a relation which we must again take into view, the fundamental relation between imagination and reality.

All reality, be it outward or inward, sensuous or supersensuous, falls --- not outside --- but inside the almighty imagination. In this very highest sense we can then say that all reality is spiritual reality; it has, after all, its source in imagination, from the creative omnipotence of spirit.

But from the standpoint of man the matter looks different. Here there is a separation, a chasm, between an outward and an inward reality, between the things outside us and our consciousness of the things. The things must be sensed; there is, to be sure, a hidden imagination in sensation, insofar as there is an intuiting; but this imagination is bound in sensation, and sensation, bound and conditioned by impressions from the things, is immediately dependent on what is outwardly given. --- This gives us an outward reality, natural reality, reality for sense and understanding. In contrast to this we have an inward world, a world of images, representations and thoughts, a supersensuous world, a world in which the supernatural too can present itself in the medium of imagination. This world has reality insofar as it has truth, and it has truth insofar as it is of divine origin. Where phenomena of this supersensuous, supernatural world come to meet us here with that clarity and power of revelation which subdues the outward, there we speak in a special and religious sense of spiritual reality, reality for spirit.

The outward natural reality impresses sensuous man and stands before consciousness with the stamp of
% --- p. 46 ---
\setcounter{footnote}{0}%
\opage{46}an independence that sets bounds to imagination and transforms its images into fancies; but spiritual man recognizes that all the outward has its origin from the inward, that natural reality with its facts and laws is and must be absolutely subordinate to spiritual reality.

From this standpoint we can now, with special regard to imagination, already see how it comes about that unbelief and superstition set themselves against faith. When the senses and sensuousness have captured and imprisoned a man's imagination to such a degree that the natural alone has validity for him, while the supernatural has been transformed into semblance and fancy, then we have unbelief. Unbelief knows no other imagination than the merely human, the aesthetically illusory. If, conversely, the supernatural has in imagination taken on the character of something outwardly sensuous, so that the spiritually real is put on an equal footing with what is material and real in nature, then we have superstition.

Only faith is able to grasp both forms of reality in their mutual difference and unity; for only faith is able to recognize an interplay between divine and human imagination. There are in the Old Testament examples in the grand style which teach us what this means.

Moses as lawgiver never speaks in his own name, but always in Jehovah's. The proclamation of the Ten Commandments begins with the well-known formula (2 Moses 20): ``And God spake all these words, saying.'' To the religious sense this cannot, taken generally, have anything particularly offensive about it. The Ten Commandments are of universally moral content; they express thoughts and principles which
% --- p. 47 ---
\setcounter{footnote}{0}%
\opage{47}every civil society everywhere must acknowledge: one could almost call them commandments of reason. The little imagination that is needed to put oneself in Moses' place, recognize his faith in Jehovah and see the necessity of his speaking in Jehovah's name when he was to speak with authority, comes as of itself. The supernatural and the natural, the supra-rational and the rational, seem here involuntarily to melt together into full harmony. But the bond of unity soon breaks when we go a little further in the text. It says (2 Moses 25): ``And Jehovah spake unto Moses, saying: Speak unto the children of Israel, that they take me an offering of every man\dots{} And this is the offering which ye shall take of them: gold and silver and brass. Dark blue and purple and crimson and fine linen and goats' hair and sealskins and rams' skins dyed red and acacia wood. Oil for the light, spices for anointing oil and for sweet incense\dots{}'' If this is verbatim from Jehovah's mouth, then Jehovah with his prescriptions for the arrangement of the tabernacle undeniably goes far into the petty; he enters into the most insignificant details. No human master builder could do it more exactly. The sublime seems here to vanish in sheer finitude. What, now, are we to think and judge from the standpoint of faith? If we take the prescriptions as dictated verbatim by Jehovah himself, then the whole passes over into outward reality and becomes an outward history. For the appropriation of such an outward history religious faith has no touchstone, no inner, unfailing mark of truth. Various religious books, the Zend books among the Persians, the Veda books among the Hindus, the Koran, and others, can, after all,
% --- p. 48 ---
\setcounter{footnote}{0}%
\opage{48}contain similar prescriptions. Credat Judæus Apella, non ego. But when the inner, unfailing mark of truth disappears, faith transforms itself into superstition. The imagination of superstition is recognizable by this, that the spiritual is sunk down into the sensuous --- goats' hair and sealskins and rams' skins dyed red can occupy imagination only in an altogether sensuous way --- so that the criterion of truth is lost. If one now goes to the opposite extreme and asserts that the special prescriptions for the tabernacle do not derive from Jehovah but from Moses alone, then one cannot, without arbitrariness, assume that the Ten Commandments are Jehovah's commandments; for with the prescription for the sabbath and the prohibition against coveting what is one's neighbor's, they too go into particulars --- the neighbor's wife, his manservant, his maidservant, his ox, his ass can, after all, also occupy imagination in a sensuous way. But when Jehovah's word shrinks to a fancy, then we have unbelief.

Where, now, is faith with its mark of truth? In the interplay between divine and human imagination; for in this is enclosed the religious-psychological interplay between the divine and the human thought, the divine and the human will, the divine and yet human vision, the divine and yet human language in divine and yet human images. The legislating word is a word of Moses and yet in truth Jehovah's word; the image of the tabernacle in all its details is a production of Moses' imagination, formed, shaped and worked out in the prophet's inner being, and yet a divine fundamental vision, the plan in each and every part Jehovah's plan. Thus speaks the God of Israel, and the speech is
% --- p. 49 ---
\setcounter{footnote}{0}%
\opage{49}real, spiritually real. Jehovah speaks to his prophet as spirit to spirit; he speaks to Moses, in Moses' inner being --- through Moses, a divine and yet human word; thus the commandments become commandments of Moses, and yet commandments of Jehovah.

What, then, do we know of this interplay? We know, in the scientific sense, nothing, nothing at all; here there is no object of knowledge, but by virtue of religious imagination there is here an object of faith. The mark of truth is unfailing when the thought of revelation is in unity with the thought of faith; the revealed will is in essence and ground the will of faith: the test of truth is in the believing mind, and the assurance in the inner being.

Let us go back to the beginning and consider the calling of the prophet. In the wilderness of Midian by Mount Horeb, far from Pharaoh's court, Moses is a shepherd. Out there among the mountains there is an infinity of solitude: memories awaken, thoughts speak, but voices are silent, silence dwells in the desolate mountains; out there there is a meeting between imagination and reality. And the real is something outward, a thornbush, and behold! there is fire in the thornbush, and the thornbush burns and yet is not consumed. Then a calling voice is heard; it speaks in Moses to Moses: ``I am the God of thy father, the God of Abraham, the God of Isaac and the God of Jacob. I have surely seen the affliction of my people. --- Go, I will send thee unto Pharaoh, to deliver my people out of Egypt.'' And Moses trembles at the thought of the great task, which is infinitely beyond his strength. ``But'' --- the voice sounds --- ``what is that thou hast in thine hand?'' A rod. ``Cast it on the ground'' --- and he cast it on the ground, and it became a serpent, and Moses fled back from the serpent.
% --- p. 50 ---
\setcounter{footnote}{0}%
\opage{50}--- ``Fancy, semblance and imagination, an outrageous breach of the laws of nature!'' These are voices from the chair of knowledge; but these voices and these cries, whether they come from the learned themselves or from their schoolboys, do not reach out to those regions, do not disturb the peace of solitude in the wilderness of Midian, nor the solemn silence that keeps watch about Mount Horeb. The only thing asked about here is the thought of revelation in its unity with the omnipotence of will and of imagination. The serpent that comes forth from Moses' rod has no place in zoology; to stare at it as at something outwardly sensuous is harmful to the eyes; one goes blind either from superstition or from unbelief. But, say those who know, a real serpent is surely meant here; Moses sees it: he shudders and starts back. We answer: yes, real it is, to be sure, but the reality is of a peculiar kind and must be seen with the eyes of faith. The supra-real is like a flash of lightning in the dark: in the same instant that it is seen, the eye has closed; no shine is so dazzling as the shine of the supernatural, no lightning is so strong as the stroke of omnipotence.

It is precisely the peculiarity of such a vision of revelation that it has the figure of something outward and resembles the natural, and yet, when sense and understanding set about investigating it, then there is either something merely natural, or the whole has vanished. Only the eye of faith can bear to look at the sign, for this eye of the soul is built and organized precisely so that the sensuous can dissolve into something supersensuous, and the outward be converted into something inward. And where, then, is the mark of truth? In the unity of thought with the omnipotence of will, in the divine imagination that makes the invisible
% --- p. 51 ---
\setcounter{footnote}{0}%
\opage{51}visible. Jehovah wills to save his people from the plagues of Egypt, from the yoke of bondage, he wills to save it; here is his word. He can save it, here is his sign! Yes, here is he himself: with this assurance, an assurance of faith which holds its ground against all the doubt of knowledge, Moses then takes up his rod and goes to Pharaoh.

That the outward natural can in spiritual reality be seized and mastered by the supernatural is unmistakable; but that in religious imagination it can also work in a free inward manner is clear from the prophetic visions. We choose a passage in Isaiah (ch. 6). The prophet describes his calling. ``In the year that King Uzziah died I saw the Lord sitting upon a high and lifted-up throne, and the skirts of his robe filled the temple. Seraphim stood above him, each with six wings; with twain each covered his face, with twain he covered his feet, and with twain he did fly. And one cried unto another: Holy, holy, holy is Jehovah Zebaoth, the whole earth is full of his glory! And the posts of the door trembled at the voice of him that cried, and the house was filled with smoke. And I said: Woe is me! I am undone; for I am a man of unclean lips, and I dwell in the midst of a people that hath unclean lips, and yet mine eyes have seen the King, Jehovah, Zebaoth. Then flew one of the seraphim unto me, and in his hand he had a glowing stone, which he had taken with the tongs from the altar. And he touched my mouth and said: Lo, this toucheth thy lips; and thine iniquity is departed from thee, and thy sin shall be atoned. And I heard the voice of the Lord, saying: Whom shall I send, and who will go for us? Then said I: Here am I, send me!'' --- If this is not
% --- p. 52 ---
\setcounter{footnote}{0}%
\opage{52}imagination of revelation, and if in this imagination of revelation there is not a true and living interplay between the divine and the human, where then is there a light of truth? If a vision with this stamp of loftiness, in this symbolism, with this purity, is not reality, spiritual reality, reality for spirit: what then is reality and what is spirit?

How religious imagination, without losing its stamp of revelation or giving up its character of faith, passes over into poetry can be seen from manifold examples in the Old Testament, notably in David. The peculiarity of this poetry is that, whether it turns inward, lamenting in sin and repentance, or turns outward and takes in impressions from the surrounding nature, admiring its greatness, exulting over its glory, everywhere imagination is filled with Jehovah's presence and sees in everything a revelation of the Creator's omnipotence. David depicts how Jehovah in the hour of danger came to his help against his enemies, and revealed himself in a storm. ``In my distress I called upon Jehovah, and he heard my voice out of his temple. Then the earth was shaken so that it trembled, and the foundations of the mountains quaked.'' --- It is in a thunderstorm that Jehovah comes. ``There went up a smoke in his nostrils, devouring fire went out of his mouth, glowing coals sparkled from him. He bowed the heavens and came down; he rode upon a cherub and did fly, he soared upon the wings of the storm.'' This presentation does not point to any special breakthrough of the supernatural; it is the human side of the imagination of faith that is here in free movement; the presentation is so clear that we can see at every point how the outward natural transforms
% --- p. 53 ---
\setcounter{footnote}{0}%
\opage{53}itself in reflections of the poet-king's religious imagination. But this imagination would not in truth be religious if it were not seized from within by Jehovah himself and inspirited by the content of revelation.

How deep-reaching the interplay between the divine and the human is in the imagination of revelation of the Old Testament shows itself from yet another side. There is an essential difference between the mythic gods, the Greek, e.g., or the Norse, and the only true God of revelation. The mythic gods are individualities of soul and body, human beings, only of a higher order. There are human inclinations and passions in these gods. They change their minds, they grow angry and act rashly; it goes so far in Homer that the gods, merely by watching the fight for Troy and taking sides, become so quarrelsome among themselves and in the end so hot-tempered that they literally come to blows till the fur flies.

Jehovah has no such limited individuality; he is infinitely exalted, unchangeable; Jehovah is spirit, his essence the eternal, omniscient, self-powerful I. When Moses asks what name he is to give him when he comes to Egypt and speaks to the children of Israel, he receives the answer: I am that I am, my name is: I shall be. And yet the Exalted One descends from his throne of eternity to the earth, the Infinite has dealings with the finite, the Invisible makes himself visible, yes, what is more, he receives a human mind, his moods change, his mercy can be turned into flaming wrath, he can repent, regret; to what extreme he can go in his exasperation when his vehemence flares up, of that we have a quite striking example in
% --- p. 54 ---
\setcounter{footnote}{0}%
\opage{54}the Second Book of Moses (ch. 32). When Aaron has cast the golden calf, Jehovah says to Moses: I have seen this people, and behold, it is a stiffnecked people. And now let me alone, that my wrath may be kindled against them and I may consume them, and I will make of thee a great people. Then Moses besought the face of Jehovah his God and said: Why is thy wrath kindled against this people, which thou hast brought forth out of the land of Egypt with great power and with a strong arm? Wherefore should the Egyptians speak thus: For their mischief did he bring them out, to slay them in the mountains and to consume them from the face of the earth; turn from thy kindling wrath and repent of the evil against thy people. Remember Abraham and Isaac and Israel, thy servants, to whom thou swarest by thine own self, to whom thou saidst: I will multiply your seed as the stars of heaven, and all this land that I have spoken of will I give unto your seed, and they shall inherit it for ever. ``And Jehovah repented of the evil which he had said that he would do unto his people.'' It is a highly characteristic situation. Unbelief exults: here one can see that faith in revelation is superstition and superstition stupid. The prophet surpasses his God: Moses is calmer, more level-headed than Jehovah himself. It is with exalted superiority that the highly trusted servant brings the agitated mind of his royal master into equilibrium. And by what means does this happen? By awakening the sense of honor of the exasperated one; for the Egyptians would, after all, laugh at him if he gave up his plan and slew those he was to set free. But above all by reminding him of his oath, so that he may not stand as a perjurer, etc.

% --- p. 55 ---
\setcounter{footnote}{0}%
\opage{55}And what, then, is it that gives offense here? The holy wrath of the spirit.

But with this holy wrath there is a peculiar connection in religious imagination. The God of revelation has a heart, a mind. He can love and he can hate, he can be angered and he can be appeased. As his thoughts, his will and the visions of his imagination pass over into the thoughts of man and set man's mind and passions astir, we see the holy wrath surge up like a human being's wrath and the divine pathos glow with a fire of human passion. We must never forget that what speaks to Moses also speaks in Moses. The mark of truth is seen when a distinction is made between faith and superstition. For superstition the whole story is outward; what Jehovah says and what Moses says stand outwardly apart, as when one human being speaks to another; then unbelief comes and shows that the whole is a myth. But faith has another explanation; it grants that Jehovah is right, unconditionally right, that the people which after such a deliverance can dance about the golden calf has deserved death; it thanks Moses as intercessor, but it sees clearly that the reference to the Egyptians' exulting and to the sworn oath does not aim at a change of state in the unchangeable essence of the Eternal, but at a movement in the inner being of the one who prays. It is the prophet's faith in the unchangeable will of the saving Jehovah that lifts him up above his own human passions. Faith understands what it means that the eternally high God of the spirit receives a human stamp, yes, a national stamp.

The God of Israel has a Jewish physiognomy; he is not always seen in the pure, clear light of heaven, but in a % print: kar (misprint)
% --- p. 56 ---
\setcounter{footnote}{0}%
\opage{56}mingled light of the Jewish popular mind, a sea of flame of Jewish passions. It is not for nothing that he calls himself the God of Abraham, the God of Isaac and the God of Jacob. With this physiognomy he is the Exalted One, his countenance mild, and yet, when wrath comes, terrible.

The imagination of revelation in the Old Testament is Jewish: rich in promises is this imagination of Jehovah, mightily overwhelming in its threats. ``I will search Jerusalem with lamps, and I will visit the men that say in themselves: Jehovah doeth neither good nor evil\dots{} It is the sound of the day of Jehovah; there the mighty man crieth aloud. A day of wrath is that day, a day of trouble and distress, of tumult and crash, of gloom and darkness, of clouds and mist; the day when the trumpet soundeth against the fenced cities and against the high towers. I will distress men, so that they shall walk like the blind, for they sin against Jehovah. (Zephan. 1).

None among all the gods has been able to bless as the God of Israel, when he blesses; and none among all the gods has been able to curse as the God of Israel, when he curses. He who would learn what religious imagination is, let him read the Old Testament. But the thought of revelation advances --- like the sun on its bright way to the height of noon; the Old Testament prepares the New; we shall see what change thereby takes place in religious imagination.

% --- p. 57 ---
\setcounter{footnote}{0}%
\lecture{Fifth Lecture.}\opage{57}

On a preliminary consideration it might almost look
as if the New Testament --- John's Revelation excepted
--- did not grant religious imagination
so widely extended or so magnificent a scope as the
Old. In a certain sense this may well be said.
For the forms of imagination become all the more prominent,
the more sharply the present, many-edged reality
draws its outlines against the bright, dawning ground of
something beyond, something unbounded and invisible, which nevertheless bounds
itself and makes itself visible. The cast shadows on
the earth become all the longer and darker, the more obliquely
the light from above falls on the jutting heights.
Light and darkness, holy and sinful, obedience and rebellion,
a stiff-necked people and a chastening Jehovah, these
extreme opposites of existence are found, at the
Old Testament stage, under ceaseless attraction
and repulsion in unceasing clashes. The Almighty
himself, the infinitely Exalted, condescends to man
and speaks with man: Jehovah speaks with Moses as
a man with his neighbour, and yet --- only Jehovah is
eternal, man is dust and ashes. Here is something that
can set religious imagination in motion.

% --- p. 58 ---
\setcounter{footnote}{0}%
\opage{58}In the New Testament religious imagination seems
to come to rest; the supernatural, in deepest
secrecy, comes together into unity with the natural,
the divine becomes human; God himself
reveals himself as man, a real human being,
who is born and grows, suffers and dies. The original
unity of God and his revelation is the Word. As
the Word itself the Invisible is visible. ``And the Word was made
flesh, and dwelt among us.'' Through the outward
lowliness shines the inward loftiness, and religious
imagination finds rest in contemplation when the Evangelist
adds: ``we beheld his glory, the glory as of the
only begotten of the Father, full of grace and truth.''
This is a manifestation that can fill religious
imagination with reality. This vision alone outweighs
and explains all the prophetic visions of the Old Testament.
``Blessed are the eyes which see the things that ye see!'' But it
must be seen in the light of faith; the eye of unbelief sees only the
offensive.

The difference between the religious imagination of the Old and of the New Testament
is grounded in the difference
between the spirit of the Old and of the New Testament.
Essentially the spirit is the same, and yet different at the two
stages of revelation. How different is
the law from the gospel!

The giving of the law on Sinai, terrible to human
imagination, is accompanied by terrors. The mountain quakes,
the lightnings flame; out of the murk and the dark clouds
Jehovah speaks to Moses and to Moses alone; no living
creature, man as little as beast, dares approach
the foot of the mountain. The effect is described in telling strokes:
``when the people perceived the thunder and the flames and the mountain,
% --- p. 59 ---
\setcounter{footnote}{0}%
\opage{59}that smoked, they fled and stood afar off. And they
said unto Moses, Speak thou with us, and we will hear: but
let not God speak with us, lest we die. --- That is
the spirit of the law in the imagination of fear; and fear is needed,
reverence for the holy, when a people still raw and
untrained is to be lifted up to the stage of spirit.

How entirely otherwise with the law of love in
the New Testament, the new law, which at one and the same time
fulfils and abolishes the old. This law does not threaten
and does not compel; it knows no curses,
no, only blessings. It is a law of freedom, a
law of the heart, for it is a law of ``beatitudes.'' We
hear it from his mouth who preaches on the mount.
``Blessed are the poor in spirit, blessed are the pure in heart,
blessed are the meek, the merciful, the peacemakers,
blessed are they which are persecuted for righteousness' sake,
--- theirs is the kingdom of heaven.'' This giving of the law strikes
no terror, this ``Mount of the Beatitudes'' is not
surrounded by terrors. Other visions, bright images present
themselves to imagination. The people do not stand afar off,
struck with dread before the Invisible; on the contrary,
they approach in confidence the Visible, they gather,
they flock around him in dense throngs. Every
word of his is a well of life, souls drink it in, the ear
drinks it like new wine. No earthquake, no
lightning and thunderclap about the mount where he preaches;
no, peace in nature and peace in the hearts of men.
How gracious, spring-fresh, solemn it is out there:
the sky high and the air clear, the birds free of all care,
and the lilies of the field so silent, still, as though
they were listening to the divine discourse.

But who is he really, the new prophet, whence
% --- p. 60 ---
\setcounter{footnote}{0}%
\opage{60}does he come, whose son is he? so it is asked all
round in the cities of Galilee, and the question sets
men's imagination in motion. It is said to be
actually so, that he is a carpenter's son; people in
Nazareth know, after all, both his father and mother, his
brothers and his sisters. --- But he speaks with an
authority as though he had come down from heaven. Whence
has he this authority? As a prophet, from
the Spirit! --- But he speaks, after all, as one who stands above
all prophets; whose son is he? --- Should this matter of
his descent remain in complete ambiguity, then he would
indeed become a mythical person; but revelation is the
opposite of myth. To the gospel of the annunciation is
then joined the gospel of the nativity.

The image which is here presented to religious
imagination has all the marks of revelation. The vision of angels
and the heavenly choir is not meant for
the shepherds alone; it is meant for all generations, that is to
say for all who, in spite of the world's deification of overmighty
geniuses of force, will turn their thought towards
Bethlehem, and contemplate the child in the manger.

In the gospel of the nativity there appears a truly peculiar
opposition and at the same time a peculiar fusion
of outward reality and spiritual reality. The child in
the manger belongs, after all, just as much as Mary, Joseph and
the shepherds, to the outward, historically real; but
the angels, where do they belong? To the world of imagination!
And yet their existence is by no means for religious
imagination what it is for unbelief, merely a superstitious
fancy or at the most a merely aesthetic-mythical
illusion. The reciprocal determination of the sensuous and
the supersensuous is here so thoroughgoing that he who
% --- p. 61 ---
\setcounter{footnote}{0}%
\opage{61}cannot see the angels cannot see the child either.
The child must be seen with the same eye of faith with which
the angels are seen; otherwise it is not the Christ child, % print: et (misprint)
but a supposititious child, not the child whom the shepherds
worship.

The two-sided reality points us towards the divine
origin of religious imagination. If
one searches in the outward for the historical sources of the gospel of the nativity,
one searches in vain; for its origin
is veiled in impenetrable darkness. In the domain of
outward reality unbelief and superstition
will be able to carry on their quarrel about it until the end of the world.

For faith the quarrel is settled; it has a mark of truth which is infallible, and that is in the spirit. Whether
Mary had given an oral account, whether the shepherds,
as many as they were, had left written testimony each one by himself
or all together: that would
not take us a single step further. The proof that
could be drawn from it would be this alone, that the shepherds and Mary had believed in the manifestation. It would
then, after all, have to become a reality of faith; but in the moment when the reality of
history in religious imagination transforms itself into
reality of faith, it is in the same instant transformed into
spiritual reality. --- ``And straightway there was with the angel
a multitude of the heavenly host praising God''\dots{}
He who, when he has fixed his gaze on the child,
cannot catch sight of the multitude of the heavenly host,
lacks religious imagination, and folios
of historical proofs will not be able to compel him
to faith. But he whose eye has once been opened
to see --- the child in the manger and the host of angels in the
holy night, for him the vision stands clear and certain: he
% --- p. 62 ---
\setcounter{footnote}{0}%
\opage{62}demands no outward testimony. He who has no
ear to hear the choir of the blessed will listen in vain
to the loud voices of criticism. But he whose inner
ear has once been opened to the heavenly message
will be able to separate truth from fancy,
revelation from myth.

When Christ had ended his Sermon on the Mount, the
people were greatly astonished at his doctrine, for ``he
taught them as one having authority.'' Whence
then has he this authority? From heaven. True
enough; but he is, after all, a real human being, born and
brought up as a human being. Authority presupposes personal
superiority. He must therefore during his preparation
have lived through something which, in a divine
and yet human way, belongs to the development of character,
and makes intelligible his infinite self-certainty, his inner personal
superiority. This leads us to
the gospel of the temptation in the wilderness.

This history of the temptation forms an organic link in
the presentation of the life of Jesus in that, in few but
great strokes, it gives us an image of the struggle in which
his resolved will reveals its purity and
strength. But that this history, as surely as it is
real, belongs under spiritual reality, reality
for spirit, surely no one can doubt. Here, if
anywhere, there is occasion to speak of the power of
imagination; for in the temptation thought and will go
indisputably together with imagination. The temptation is
no selfish sensuous ferment in Christ's mind, but a
trial of strength with the spirit of darkness. Here is a collision
of will with will, Christ's will and the devil's
will.

% --- p. 63 ---
\setcounter{footnote}{0}%
\opage{63}The two wills keep themselves apart, but the real
separation, which conditions the collision, is not in
outward space; it is in the medium of imagination. The outward
and merely natural, the long fasting, the solitude of the
wilderness, the desolate surroundings, all such things form here
a relief for the inward.

There are three thoughts of will, all of which concern misuse
of the supernatural gifts with which a Messiah, the Son of
God, must be endowed, or rather three forms of
the same thought of will, which the tempter seizes and applies.
The first is a misuse of the gift of miracles,
which aims at the immediately sensuous: ``If thou be
the Son of God, command that these stones be made bread''. That
was an imagination which, humanly seen, could tempt
the hungry one. Had Christ given way to this
temptation, then by that very act his mission would have
been in vain. He who in sensuous selfishness will use a
supernatural gift to still his hunger, his thirst,
will also in sensuous selfishness use it to avert
sufferings, insults, and to defy the mockers at Golgotha
by coming down from the cross. It belongs to the
sacrifice of love which he who consecrated himself to renunciation
brought, that he who by a miracle could feed thousands
in the wilderness would not touch a stone in order by a word of omnipotence
to satisfy his own hunger. ``Then
the devil taketh him up into the holy city, and setteth
him on a pinnacle of the temple''. Does this event
gain in real understanding by faith's transforming
itself into superstition and making the whole into an outward history? % print: sil (misprint)
Christ would then have to have been carried through the air by the devil
roughly like a Doctor Faust on Mephistopheles'
cloak. Here nothing is gained, but here everything is lost: unbelief
% --- p. 64 ---
\setcounter{footnote}{0}%
\opage{64}is proved right when it transforms the whole into a myth.
It is in the spirit, borne by a demonic imagination,
that Christ stands on the pinnacle of the temple. The suggestion
to cast himself down, so as to let himself be borne up in the hands
of angels and transform his messianic appearance into
a grandiose self-glorification, a divine adventure
in which there would be need of neither patience nor obedience
nor the pain of self-abasement which was to end
with a death on the cross --- this suggestion comes
not from Christ's mind, no, it comes from the being
whose will is rebellion against God. The Messiah lets himself be lifted
up to this height by the spirit of darkness; for he does not grow dizzy
on the pinnacle of the temple: the tempter's attempt is to him a
means to fulfil the earnestness of his will, to consecrate himself to
unconditional submission to the determined counsel of his heavenly
Father.

That the struggle is between spirit and spirit, the spirit of darkness and of
light, of falsehood and of truth, a struggle of will in
the medium of imagination, becomes especially clear from what
follows. ``Again, the devil taketh him up into an exceeding
high mountain, and sheweth him all the kingdoms of the world, and the
glory of them''. Where, then, lies the high mountain from which
there is a view over all the kingdoms of the world and their glory?
In the land of imagination, on the coast by the open
sea --- in the deep dwells expectation, and the waves of possibilities
roll --- from there is the view: on the wild current
out yonder lies the isle of the Sirens, storm-beaten ships sink
as wrecks; high towers are seen in the air, palaces and castles
in the bank of clouds, great cities and far-stretching regions which
conjuring fancy peoples. On the coast yonder, in
the land of imagination, there lies the mountain, the high mountain
with a view over all the glory of the world.

% --- p. 65 ---
\setcounter{footnote}{0}%
\opage{65}Many a one, highly gifted, richly endowed, many a
hero whose name still glitters in the annals of history, has
stood on the high mountain with the possibilities before him
and the conditions of the world around him; but he mistook
semblance for reality. Spiritual reality was to him only
an imagination; then came the tempting will, whispering,
creeping forth from its hiding place, veiled in figures as in a myth.
The chosen one chose the outwardly great; egoism struck
the strings of reality --- a powerful stroke! The strings
sounded, roared and burst, but the fundamental chord rang on
misused spirit: the chosen one fell and drew millions
after him in the fall. How he rose, how high he
stood and how deep he fell! History looks at the
real, and the gigantic reality is like a
myth. So it is, then, for the salvation of the world that One has
stood on the high mountain, a single One, who with full
consciousness met the spirit of darkness and saw through the deception.
So it is, then, for the salvation of the world that, in spite of
the loud cries for strength of will and greatness of will in the beguiling imagination
of self-deification, there can still be heard a
voice, an authoritative voice, which sounds from the mountain:
``Get thee hence, Satan: for it is written, Thou shalt worship
the Lord thy God, and him only shalt thou serve.''

Is what has been set forth here now really fitted
to convince every reasonable human being of the truth of revelation,
particularly of the truth of Christianity?
We answer openly: No, by no means! Were the truth of revelation
to be grounded with universally valid, incontrovertible
proofs of reason, then it would indeed have to be grounded
scientifically. But no educated person at present will find
it objectionable that every attempt to get either
% --- p. 66 ---
\setcounter{footnote}{0}%
\opage{66}the history of the annunciation or of the nativity or of the temptation
secured as real history, whether by historical or by general philosophical
grounds, deserves to be shown out of
the precincts of scientific justifications. With the whole content
of revelation we find ourselves by no means
in the sphere of knowledge, but in that of faith. He who declares that
faith and knowledge of faith have no validity for him,
that everything about the supernatural is against his nature and
acts as a pressure on the pit of the stomach, a prelude to something
that begins with nausea, has not only the right to
declare himself, but also a claim to a certain recognition,
in so far as he speaks his mind plainly
and without circumlocution announces himself as an unbeliever.

It is otherwise, on the other hand, with those
who in the name of knowledge and science attack revealed
religion. They appeal not to their unbelief,
but to their love of truth. If it could be seen in the
clear, infallible light of knowledge that there was
truth in faith in the supernatural, then they would
gladly believe; but it is so improbable, and they are so
truth-loving! It is either a self-deception or a
hypocrisy. Science ought to have its due; but it
may not overstep its boundary. It is necessary
that the decisive point of contention be kept sharply
in view. Here on the boundary between faith and knowledge
the question is not of probability; here the question is solely of
possibility. Let science demonstrate that a phenomenon of revelation is improbable, highly improbable,
yes, the most improbable of all: that says nothing, nothing at all!
But if its impossibility can be proved, then the sphere of faith
is broken through, and its light is extinguished. Modern
science keeps watch at the boundary, and every
% --- p. 67 ---
\setcounter{footnote}{0}%
\opage{67}time it makes a sally, the battle-cry is heard: ``Impossible!''
Let us briefly see whether this ``Impossible'' should not
be a false imagination.

We cast a glance at the doctrine of light. Physics
has set up the highly probable hypothesis that light
rests on a vibration in the aether. He who acquaints himself
with higher optics will in this respect too
have rich occasion to admire human
acuteness. It is proved that red light
must give over 500, violet over 700 trillion
vibrations in a second. But is the essence of light
hereby explained? By no means. Light as light, as
clarity, presupposes a sight, a seeing consciousness, which
cannot be derived from vibrations of the aether. ``When the aether
is at rest, there is no light''; but when all eyes are closed,
when there is no seeing being, either in heaven
or on earth, then there is no light either. To speak
of a world-light which is independent of the eyes of all individuals
is to speak of a world-sight which is independent
of all seeing souls, that is, to speak of the Creator's eternal
seeing, of the divine imagination of omnipotence.

If space and the aether and the conditions
that set it in vibration did not originally proceed from
the divine seeing, then the doctrine of light would have to
be a doctrine of an inexplicable semblance without essence.
Is there now in the whole scientific doctrine of light a single
point of support for the assertion that the creative word: let there be light!
should belong to the impossible? The word is indeed precisely a
categorical expression for the unity of the divine thought
with the divine will in the all-seeing power of
imagination. In the form of imagination the word is the
original expression of essence for the procession of the visible
% --- p. 68 ---
\setcounter{footnote}{0}%
\opage{68}from the invisible. Can there now be shown in space, in the aether,
its waves and vibrations, a single fundamental determination
which proves the impossibility of the infinite seeing of an all-seeing
Godhead? None, none at all. ``Impossible''
is here an empty outburst of empty fancy, a false imagination.
But if science cannot prove the impossibility of the creative word
when the talk is of the Invisible's
making the natural universe visible, where then will it
go with its ``Impossible'' when the talk is of the free
revelation of love in the word of personality itself, the
Word which is ``the light of the world,'' the creating Word which
was in the beginning?

Objections which for a scientific criticism
are of greater significance are drawn from the conspicuous
opposition between what happened in Christ's time, and % print: Chisti (misprint)
what in our time is taken to be possible. ``Who can in
our day believe in possessions by devils or raisings
of the dead? Even if, as far as possessions by devils are concerned,
one may take refuge in the shadowy hiding-places of religious imagination
and there seek cover from the pursuing
criticism, it will nevertheless not do to make the raising of the dead
into an imagination. To raise a dead man is
just as impossible as to make a stone into a human being.
He who believes in such things must break the staff over all
the results of science. One of two things: either one must
put out the light of reason and give oneself up to the darkest
superstition, or grant science, with its inductions,
analogies and far-reaching experiences, the right to
demand that the principle of an uninterruptedly coherent,
evenly progressing natural development be respected, a
principle which is of all the more weight as
it alone accords with what reason teaches us about
% --- p. 69 ---
\setcounter{footnote}{0}%
\opage{69}the immutability of the divine being.'' --- As long as
the standpoint remains standing within the boundary of the probable
and the improbable --- hence outside the circle of light of the knowledge of faith
--- such objections will make
an impression; but when the question is of possible and
impossible, they dissolve into a nothing. When science
maintains that in nature it is not able to discover
anything that could point towards a supernatural, and
thereby sufficiently demonstrates that the doctrine of the supernatural
does not belong under what can be scientifically
grounded: has it then by the same token proved
that the intervention of the supernatural in the course of nature is
impossible? When science maintains the immutability of the divine
being, has it then by the same token
proved that this immutability must necessarily make
the divine self-determination, the freedom of God's will,
impossible? When science maintains the impossibility that
the absolute spirit can contradict itself, has it then
by the same token proved the impossibility that this
natural order of things can, in its invisible connection, accessible to
human knowledge, with the
supernatural, be a transitional link and means for
higher purposes? Has one perhaps proved the impossibility of such a % print: maaksee (misprint)
higher purpose, or the impossibility that
a means can serve its purpose? When science
maintains --- and it can be maintained even without science
--- that a dead body cannot be called to life by
natural means, has it then by the same token proved
the impossibility that it can happen by supernatural
means? --- Unbelief thinks to make itself unrecognizable by
wrapping itself in the cloak of knowledge, but the wind blows
in the cloak, and by the claws one knows the lion.

% --- p. 70 ---
\setcounter{footnote}{0}%
\opage{70}Here there is, be it noted, by no means talk of essential
discord between knowledge and faith. Negative criticism must
be allowed to prove that the faith of revelation is without
scientific validity; if only it proves it properly
scientifically, purely scientifically, then one knows what
it is: it does no harm at all. Here there is talk only
of the denial of faith which, under the semblance of love of truth
and of strict scientific rigour, has grown
self-important in the endeavour to undermine
the religious foundation of religion.

Unbelief is enterprising, shrewd and politic; it understands
so well, where it suits its purposes, how to ally itself
with superstition, reject spiritual reality and demand the
outwardly tangible. --- ``If there is,'' says negative
criticism, ``anything which faith must acknowledge as real,
outwardly historically real, then it is surely Christ's
resurrection.'' Yes! answers superstition. ``Does not Christ
after his resurrection go about with the very same
material body as that with which he was buried?''
Yes! answers superstition, with the very same material
body. ``Is it not veiled evasion and half-faith
here to talk of spiritual reality and resort to the help of
imagination, in order to escape the objections of doubt?''
Yes indeed, it is veiled half-faith, this spiritual reality,
answers superstition. --- When unbelief has now by
this road got superstition onto slippery ice, it gloats
and turns the whole thing inside out: ``Christ makes himself
invisible, he vanishes from the two who went to Emmaus;
therefore he has no material body. Christ goes
through closed doors; therefore he has no material
body. Christ ascends to heaven, contrary to the law of
gravity; therefore he has no material body.'' --- The
% --- p. 71 ---
\setcounter{footnote}{0}%
\opage{71}out-and-out historical superstition stands nonplussed, suspects that it has
drunk brotherhood with the executioner, and knows neither which way
to turn. Unbelief smiles, chucks the nonplussed one under
the chin, looks it in the eyes and says: Reality!
Yes, you are after all a good honest piece of reality; he
who says of you that you are a spiritual reality, reality
for spirit, lies shamefully about you. Don't
be angry; at bottom I quite like you; but ---
you are stupid.

Faith itself unbelief cannot abide; their
mind and vision are infinitely different. It is not only
reality, but also imagination, about which they
disagree. All reality of faith is, even if it is not
expressed with clear consciousness, always spiritual reality;
but spiritual reality is an entrenchment which the
negative criticism, with all its shrewdness, its apparatus of reflection
and tactical manoeuvres, is not able to
break through. When it thinks it has got hold of
the outward, then this transforms itself and passes over into an
inward, and when it then thinks it can hold fast
the inward for itself as imagination and fancy,
then this transforms itself again and breaks out into an outward.
Impatient and embittered, criticism bursts out: ``let
us be told what and how much here is outward,
let us be told what and how much is to be reckoned
to the inward!'' But faith answers unperturbed:
here is the mystery, and the mystery is no task for
knowledge. Unbelief fancies that it strikes the decisive blow
and hits the nail on the head when it sums up
the result of the investigations in deterrent
catchwords: ``Mystique, mysticism!'' Nor does this make
an impression on faith. The mystery does not tempt to
% --- p. 72 ---
\setcounter{footnote}{0}%
\opage{72}overstrained intuition or fruitless fantasizing; it indicates
briefly, categorically: the boundary of human knowledge.
The mystery of revelation gives no direction to
erase the marks of reality and to sneak round
the laws; on the contrary it gives direction to a law, a
fundamental law of freedom, which is just as certain for faith
as the laws of science are for science; this
fundamental law is a law of transformation.

The law of transformation comes to view in the wonderful
working of the supernatural upon the natural. That this
law is not taken in vain is secured by its
being a law for the improbable. Faith has for its
part just as sharp an eye for the improbable as
unbelief for its own. If the eye for the improbable is dulled,
faith passes over instantly into superstition; for superstition
is stupid. But faith is not stupid; it knows what it
is doing when, in spite of the improbability, it seizes the possibility.
Herein lies this: that the possibility which faith seizes against all
probability must, for the salvation of the soul, for the
peace of the heart, be of such thoroughgoing significance that
it cannot be given up without the assurance of
eternal life being given up by the same token. Faith too
knows what it is doing: it has its reasons.

We can test the law of transformation on the gospel of the resurrection.
Surely --- that is the assurance of faith
--- Christ is risen from the dead with his
real body, not with a semblance-body. But that this
real body nevertheless rises above the merely
outward and material, that it already belongs to the sphere
of spiritual reality, makes itself plainly known
in the words of the Risen One to Mary Magdalene in the garden.
The supernatural works through the natural.
% --- p. 73 ---
\setcounter{footnote}{0}%
\opage{73}The body of the Risen One undergoes a series of transformations.
These transformations, which correspond to the various
manifestations to the disciples and end with
the ascension as the last transfiguration, belong essentially
to the conclusion of the life of Jesus, because in his
personality they make manifest the original fundamental relation between
the natural and the supernatural.

The law of transformation in and for itself is so far from
standing in conflict with the great principles of natural science that,
on the contrary, it is this law alone which is able to ground
the right of reason. To understand nature without spirit
from bare matter is impossible. The philosophy of nature
leads us to the incontrovertible result that nature
is created, borne and explained by spirit. The boundary between
faith and knowledge is determined by this, that knowledge must hold
its result in the form of necessity, and cannot make
use of the possibilities of freedom which dawn therein. It
is faith's unassailable presupposition that the fixed,
present order of things, however evenly it develops,
however many millions or trillions of years its
course of development displays, nevertheless has its innermost
original ground in a disposition of creative freedom; that,
in other words, the whole universe rests on the foundation
of a hidden wonder.

Revelation makes the hidden manifest. The miracle of the resurrection
makes intuitive the relation between the
corruptible and the incorruptible, between here and
hereafter, between death and life. Of the mystery of revelation
physics with all its mathematics knows nothing,
nothing at all. It has in ``its auditorium'' the right to speak
of the specific gravity of water and of bodies, of the universal validity
of the law of gravity in the natural world, and the like; but
% --- p. 74 ---
\setcounter{footnote}{0}%
\opage{74}when unbelief on that account makes an errand into
the circle of light of revelation, in order to cry in the name of the law of gravity
and of specific gravity: Impossible! it betrays a spiritual crudity
that must put up with being shown outside.

It shall be readily conceded that the content of revealed
religion has yet many other sides
from which it must be illuminated besides precisely that of imagination. Nevertheless
it may well be time to become aware
that the great confusion in the domain of faith derives in
large part from a stunting of religious
imagination. While in all circles of culture
care is diligently taken that the aesthetic imagination, with
sense and taste for everything that belongs to art
and poetry, is duly developed, religious imagination
has stepped into the shade, even to such a degree that its
education seems to be concluded with the age of confirmation.
A mixture of all sorts of elements, a little from
aesthetics and a little from physics, a little from mythology and a little
from scientific criticism, confuses the vision and obscures
the apprehension of the word of revelation in the New as well as
in the Old Testament. But imagination is a power,
in religion too; everyone who feels a need to
give himself an account of his faith will surely also feel
a need to clarify for himself his religious imagination.

% --- p. 75 ---
\setcounter{footnote}{0}%
\lecture{Sixth Lecture.}\opage{75}

What is limited in one respect can be unlimited in another respect; this evidently holds of the infinite realm of imagination. That this realm is infinite must be granted even by those who deny its divine nature and seek the origin of imagination in a physiological disposition of nerves and brain.

In every activity belonging to the life of consciousness there is a hidden imagination. No memory is possible without imagination; the images of experience, images of what we have lived through, of what we have seen in nature or learned in history, do not lie stored away in the soul as in a warehouse: they are within us as possibilities, as dispositions in our nervous system. That we remember what is past rests on this, that imagination, at a prompting of the will, can in an arbitrary manner awaken the dormant possibilities and set forth what we have lived through in an image before our inner sight. The understanding too, with its sober, law-bound activity, is conditioned by imagination. It is the understanding that separates, compares, orders and connects our representations; but what is a representation? A dissolving image, hovering between the invisible and the visible, the indeterminate and yet determinate: as the dissolution of the image is, so must
% --- p. 76 ---
\setcounter{footnote}{0}%
\opage{76}the image also be, but the image is formed by imagination. Mathematics is in the strictest sense a science of the understanding; but mathematics is impossible without a corresponding imagination. The countless combinations of forms and formulas, the manifold compound, doubled, often highly intricate geometrical figures make it evident enough that there must be a mathematical imagination of its own. Nor can the will, however practical it is and however exclusively it aims at the real, do without the assistance of imagination. The task that the will sets itself is, after all, not yet really solved, but it is to be solved, that is to say, it is a possibility whose content must be intuitable in the form of imagination. And how intimately is imagination not united with feeling; think of mood, the mood of joy, the mood of sorrow: no ideal mood without imagination!

But imagination, however much it is at the service of every activity of the soul on every occasion, is nevertheless not content with a simple role as servant; it has a secret, one might say feminine, drive to rule itself. It does not go about it clumsily; it proceeds warily for the sake of the enchantment: it uses cunning. It does not announce itself as a loud-mouthed master who stamps on the floor with clinking spurs and gives the word of command; no, as a nimble little girl, with a character reference: deft at her work and willing for anything.'' But precisely on account of this deftness and willingness one must keep an eye on it, constantly have it in the spyglass and closely watch its every step. It looks so innocent and is yet so calculating, cunning, scheming, an actress of the first rank, with mastery in transformations, magic pageants and bewitching illusions.
% --- p. 77 ---
\setcounter{footnote}{0}%
\opage{77}Imagination's repertoire of roles is of an enormous range, from the most profound and moving to the merriest, most wanton, foolish and silly that there is.

Although imagination is sober, pliant, mobile, and can take up everything and work over what it takes up in the freest manner, it nevertheless has a different stamp in the different individualities. One sees this in what is peculiar to every kind of talent. There is a painter's imagination of its own, a musical imagination, the poet's imagination, etc. Where the gift of imagination in one direction or another oversteps its limit, the strained and mannered begins at once. Indeed, not only in the specially gifted but in every well-organized individual, imagination has its own fundamental form. This is understandable enough when we consider that the unity of intuition in which all impressions of life, thoughts, representations, models and endeavors are taken up is supplied precisely by imagination. Just as every human being in the bodily sense has his peculiar face, so he also has in the sense of the soul his peculiar physiognomy, determined by his fundamental imagination: imagination is the ideal face of the soul. Without a higher reciprocal semblance of the two faces, the inner and the outer, the mimic element in acting would lose its significance. To be a great mime takes something more than being able to set one's face in folds and pull grimaces. All plastic inspiration passes through imagination.

This is all the more natural as imagination also comprises that side of consciousness which through the half-conscious reaches down into the dark regions of the unconscious. Imagination works just as well in the night of consciousness as in its daylight; it loves the twilight,
% --- p. 78 ---
\setcounter{footnote}{0}%
\opage{78}the dusk and the waking dream, in which its airy plays are most gladly performed. In keeping with this, the power of imagination does not depend immediately on conscious thoughts or an arbitrary course of representations, but permeates flesh and blood and works as a spiritual force of nature. Think of the anxiety and terror that a glance or a word can awaken; think of the involuntary blanching and the equally involuntary blushing.

Therefore one should never measure a human being's personal cultivation by the development of his understanding, his insights and his attainments alone; if we would have a measure of his strength of mind, we must take imagination into account as well, the elemental imagination that dwells in the nerves. If this imagination is warped or neglected, the will has no secure foundation, and self-possession can fail one at the decisive moment. Composure, collectedness, presence of mind rest essentially on immediate imagination. A coarse human being, whose imagination is as blunted as his conscience, can in certain cases show greater cold-bloodedness than a cultivated man whose elemental imagination has not been developed in proportion to the enlightenment of his understanding. A church thief who breaks in by night can grope about by the altar and among the coffins in a burial vault with a calm as untroubled as if he were walking in his own sitting room. If someone suddenly seizes him from behind in the dark, it still does not occur to him that it is a ghost; he knows at once what it is: it is the police. On the other hand, many an enlightened man can hold himself fully convinced that there are no revenants, and yet by night be afraid to seek the company of the dead: he is not really afraid of the dead, but he is afraid of his imagination.
% --- p. 79 ---
\setcounter{footnote}{0}%
\opage{79}The well-known La Mettrie, one of the very crassest materialists toward the end of the last century, a man who was afraid of neither God nor the devil, was afraid to walk across a churchyard at night.

He who has a sick imagination cannot possibly have a healthy soul. As one speaks of hardening, e.g.\ in the liver, so one can also speak of knots of hardening in the imagination; what else are the so-called fixed ideas but such knots? All diseases of the soul, whatever their organic causes may be, point to a disturbance in the imagination. The understanding is, to be sure, also disturbed; but that the confusion of the understanding comes from disturbing images, wild promptings, can be seen among other things from the fact that the mad can reason with sharp understanding and sure logic within the limits that the fixed ideas mark off. How uncanny to take a look into a madhouse: all the distortions of imagination become visible there.

There is a picture by Kaulbach; it represents a madhouse. One sees the insane gathered and grouped, but one also sees that the one takes no notice whatever of the other: each of them is sufficient unto himself, exclusively occupied with his own fixed idea. The most varied fancies and passions stand painted in faces and postures. There is one in particular --- the figure is no caricature --- but he arouses a shudder. It is an enthusiast who is proclaiming the gospel; he lifts up the cross and smiles like one transfigured. But in this look, in this blissful smile, what a bottomless abyss of madness!

Imagination is an essential form for the content of truth, and yet the first inventor of every kind of fancy, lie and deception. From the psychological side there lies
% --- p. 80 ---
\setcounter{footnote}{0}%
\opage{80}in this a summons to examine more closely the origin of imagination and to give an account of where it gets its highly heterogeneous content. Psychology in our day must lean on physiology; without insight into the structure of the organs and the organic conditions it is impossible to arrive at clear knowledge of, or fruitful insight into, the activities of the soul. In this respect too the advancing science of nature has gained a considerable influence on the science of spirit. Here belongs a closer elucidation of the so-called hallucinations. It is well known how the fever patient can fantasize, how the drunkard in his delirium gets visions of so palpable a reality that he grapples with them as with bodily objects. The opium eater's imagined splendors, and their gradual transformation into just as many real miseries, belong to the phenomena that attract attention.

In J. L. Heiberg's collected writings (8th volume) there is a little essay on the Confessions of an English Opium-Eater.'' ``Opium,'' it says there, ``brings the soul into a calm, indeed serious, state of enjoyment. All the faculties of the soul seem to develop to an extraordinary degree, while at the same time they maintain the most perfect harmony.'' We get here examples of what a material influence on the nervous system has to signify. ``I stood,'' says the confessor, ``high above the tumult of life; confusion, care, anxiety, the secret burdens that oppress the heart, all had vanished; it seemed to me like a holy eve for the whole of creation, like a day of rest from all the miseries of mankind. The joys and hopes of earthly life were reconciled with the grave, and my spiritual motion went on without
% --- p. 81 ---
\setcounter{footnote}{0}%
\opage{81}exertion, without weariness, like that of the stars in the heavens. And with all this a deep calm, which did not arise from a state of inactivity but seemed to be the result of the opposition of infinite forces; it was at once an activity without limit and an infinite rest''. With what magic a morbidly strained imagination can work, to what degree it is able to conjure up a world of visions and give them the stamp of reality, shines forth from the following description: ``Transported beneath the sky of the tropics, I saw myself'' --- one should remember that the man is in London --- ``now surrounded by all the extraordinary productions, both vegetable and animal, of these regions. Birds, monkeys, reptiles surrounded me on all sides and pursued me incessantly. I fled into pagodas and found myself shut up forever in the secret recesses of these buildings. I was at once the idol and the priest, the god and the sacrifice. I fled through the forests of Asia, pursued by the wrath of Brahma; Vishnu hated me, and Siva laid snares for me. All at once I stood face to face with Isis and Osiris; they accused me of having committed a crime at which the ibis and the crocodile trembled and quaked. They condemned me to a thousand-year night in sepulchers of stone, and buried me with sphinxes and mummies in inaccessible pyramids. Then I found myself sunk in the mud of the Nile, in the midst of monsters, steeped in river slime, and among all kinds of horrible and loathsome things''. --- The images grow ever darker and wilder. ``A tremendous war, a battle, a death-struggle was to take place somewhere or other, and unroll itself before me like a dark and mournful drama; but I knew neither where nor
% --- p. 82 ---
\setcounter{footnote}{0}%
\opage{82}how, nor by whom. The lively concern I felt for it was all the more unbearable as I knew neither the nature nor the cause nor the probable outcome of that in which I took an interest. As often happens to us in dreams, where we are always the center of all that goes on around us, so it seemed to me here too that I had the power to determine this outcome, and yet at the same time that I did not have it. I felt that I could do it if I had the strength to will it, but I was as if weighed down by a tremendous burden from which I could not free myself. I therefore remained in my inactivity. But now the critical moment drew near, the tumult grew stronger, an infinite matter was to be decided. Then followed a sudden consternation, headlong flight in all directions, and the echo of the footsteps of countless fugitives. Then shadows and torches, whirlwinds and human faces passed me by; and finally --- with the feeling that all was lost --- female forms, and the features of what had been dearest to me in the world. And only a single moment was granted us; we had to part forever. All wrung their hands, lamented and howled, and from all sides one heard it repeated: Farewell forever! --- These dreadful words echoed within me; I awoke in the most terrible anguish and cried with a loud voice: I will never sleep again!''

That such things must come from affections of the brain, there can be no doubt. Physiology has in this respect great merit for having given an account of the different kinds of hallucinations.

Yet in this point too the science of spirit must be at its post, for the matter can easily take a turn which
% --- p. 83 ---
\setcounter{footnote}{0}%
\opage{83}distorts the conception, reduces the ideal to semblance and perverts the very highest.

Just as criticism, with the help of a deeper-going psychology, has discovered the origin of myths and transformed revelation into myth, so it can now, seen superficially, with a certain right appeal to a deeper-going physiology and dismiss the historical testimonies that bear the stamp of psychological credibility, where the talk is not of obscure, half-conscious, but of clear, fully conscious manifestations, with the short, dry remark: these are hallucinations.

A couple of examples. Saint Anthony flees from the world out into the desert, but the devil comes along and torments him with all manner of temptations. Now it is hideous trolls that terrify him, and now lovely women who put his abstinence to the test; his cries in the cave are heard far off, for Satan himself attacks him bodily and strikes him to the ground. In this, says the physiologist, there is nothing unreasonable. It may all really have come about so. The hermit has tormented himself with ascesis, fasting and mortifications for so long that he has become a victim of the most terrible hallucinations.

It is told of Mohammed that when, in solitude, with prayer and fasting, he brooded over his new religion, he heard from all sides in the surrounding nature voices that confirmed his divine mission. It spoke from the rocks: Thou art the Prophet; it sounded from the clouds: Thou art the Prophet; the trees put their heads together, the foliage whispered and the springs murmured: See, there goes the Prophet. And Mohammed was led through all seven heavens. In the highest he saw the angel of death;
% --- p. 84 ---
\setcounter{footnote}{0}%
\opage{84}indeed, he saw God's greatest creature, an angel that had 70,000 heads, in each head 70,000 mouths, in each mouth 70,000 tongues; with each tongue it spoke 70,000 languages and proclaimed in all of them the praise of the Lord. At last he was brought before the Lord God himself, and to the right of the throne there stood to be read: there is no God but God, and Mohammed is his Prophet. When the tribe of the Koreishites set themselves against him and would not listen to his preaching, because he infringed their prerogative by fighting against idolatry, his paternal uncle Abu Taleb was the only one who took him under his defense, although he refused to profess his faith. He advised his nephew to desist from his undertaking, since otherwise great misfortunes would flow from it; but Mohammed answered: ``Though they came to me with the sun in the one hand and the moon in the other, they should still not make me yield''.

What, now, is fancy and what is truth? Can this be decided in the sphere of faith with objectively scientific grounds and proofs? Impossible! Even if we had in every city 70,000 theological faculties, and every faculty 70,000 learned theologians, and every theologian 70,000 folios, and every folio 70,000 theological proofs: it would still not shake a Mohammed's faith, would not, even if he had understood every single proof to the letter, make him yield a hair's breadth. Can this then perhaps make us uncertain of the validity of the knowledge of faith? Not at all. The sense of sight can deceive; why should imagination not be able to deceive too? But as little as any reasonable human being will give up trusting his sight, although he knows that sight can mislead, just as little will any believing soul
% --- p. 85 ---
\setcounter{footnote}{0}%
\opage{85}take offense that revelation comes in the form of imagination, although he knows that imagination can mislead.

In order to distinguish sharply between true and false, we must constantly keep our attention fixed on the essential content of imagination in its unity with the form. This holds not only of the religious, whose content is spiritual reality, reality for spirit; it holds also in spheres where the content is what we specifically call idea. It holds in art and poetry; it holds practically in rational ethics. Art has its truth, although there are distortions of art; poetry has its truth, and this truth is recognizable precisely by the fact that there is an untrue poetry, an empty and insipid, indeed immoral poetry, and one highly corrupting for the unsteady.

May poetry be immoral? There has been much dispute about the relation between the beautiful and the good. Those who demand that poetry be moral in the sense that it should immediately moralize forget that, as regards the form of content, there are two essentially different ideas, the aesthetic idea, the beautiful, and the ethical idea, the good; they forget that imagination is sovereign in the realm of beauty and has the right to take everything into its service, not only the good but also the evil, the immoral, the hideous. But those who go to the opposite extreme and assert that poetry need pay no heed at all to moral demands, those who assume that free imagination, like free thought, can in the ethical sense offer the reader whatever it pleases; that, with reality and its possibilities in view, it celebrates its most brilliant triumph precisely when it knows how to seize upon a brothel story,
% --- p. 86 ---
\setcounter{footnote}{0}%
\opage{86}a subject in its prosaic outlines so vulgar that mature men, for decency's sake, push it from them with distaste, and yet to furnish this shabby subject so splendidly, so brilliantly and piquantly, and to let the magic of the presentation glorify the indecent even to such a degree that young women readers, without knowing what they do, become absorbed in it and drop by drop have their imagination poisoned; --- those who conceive the poetic from this standpoint must have shut their eyes beforehand to the divine mission of poetry.

It is not merely what one may perhaps, from a high-flying standpoint of emancipation, call a philistine morality that is thereby mocked: it is the idea itself that is violated. Make the gulf between the aesthetic and the ethical as wide as you will, it is surely still evident that the ethical idea has a reality at least as valid as the aesthetic.

That there must be an ethical as well as an aesthetic imagination is easy to make clear. But in the ethical, imagination serves. It has already been remarked that the will needs the help of imagination. Nothing is more dangerous for the moral will than a confused, bewitched imagination, led astray by false images. Whence comes vanity? From false imagination! And ambition and avarice? From bewitched imagination! And the demon of lust? From poisoned imagination! He who systematically sets out to corrupt a human being will not confine himself to overturning these or those principles, or to luring the person in question into these or those habits; no! The corrupter cannot be sure of his victim until he has thoroughly corrupted its imagination. Only with the
% --- p. 87 ---
\setcounter{footnote}{0}%
\opage{87}corrupted imagination would the last stirrings of the pangs of conscience cease, and the last remnant of a moral idea be torn up by the roots.

But it is surely evident that the true ideas, however different they may be among themselves, cannot possibly be each other's enemies and harbor such contempt, indeed so ingrained a hatred, for one another that the one should set out to scorn or even to exterminate the other. The relation is, after all, just the reverse: there is at bottom good understanding between idea and idea; there is peace and harmony in the sphere of ideas.

One can, to be sure, also speak of an egoism of ideas, inasmuch as a single idea can press forward so strongly that it pushes the others aside; but this egoism does not belong to the inner nature of the ideas; it comes with their passage through the human mind, at the transition from eternity to time, to emergence into the world, to realization in the temporal. It is with the idea in man's mind as with the yolk of the bird's egg: first it is small and soft as a berry, next like a grape, and at last, in its passage through the oviduct, it is coated with the membrane and the hard shell. And yet the idea comes forth from its germ otherwise than the bird from the egg. The more an idea is inspirited by the universal and permeated by spirit, the more clearly will it show us, in its individualizing appearance, that with all its peculiarity it acknowledges its kinship with the other ideas, especially with those of the nearest circle to which it belongs. This holds literally of the relation between the aesthetic and the ethical idea.

``Modern poetry'' does well to take ``the children of the world'' as they are, and to lay bare the hollowness and conceit
% --- p. 88 ---
\setcounter{footnote}{0}%
\opage{88}with which a bleached imagination of piety or a regimented barracks morality is so fond of talking big; it does well to reinstate the drives and dispositions of natural man in their original rights. But these drives and dispositions must, after all, have a human goal and be subject to the laws of spirit. It is an uncanny sign of confused imagination when not even the poet can catch sight of the laws of spirit, but makes the satisfaction of drive itself into an ideal, makes his heroes and heroines into meteors and soap bubbles, and smiles ironically when the strong souls, entangled in their own consequences, go under. One cannot reproach modern poetry for immersing itself in the real and gaining in mastery where it is a matter of assimilating its elements in living forms and situations pregnant with character. But one can with reason reproach it that its images of life are false, because they have either an insipid and empty or even an altogether false background. The poetic does not suffer from the present, real life, with its contrasts of self-important narrow-mindedness and ecstatic emancipations, being allowed to step forth in clear features; but it suffers terribly when the impression of the whole presentation shows that the poet himself, with all his gifts, is impoverished in spirit and merely plays with ideas, a worldly-wise and yet, in sensuous emancipation, confused subject who lacks a standard. Modern poetry --- Cherbuliez's Miss Rowel, e.g.\ --- understands very well how to reconcile the ethical with the aesthetic, insofar as it has any of it, but it does not have much. % print: ekstastiske (misprint)

In their eternal essence the good must absolutely be beautiful, and the beautiful good. Nonetheless the
% --- p. 89 ---
\setcounter{footnote}{0}%
\opage{89}form of imagination in which these ideas appear is highly different.

In art and poetry imagination works over the real solely from the side of possibility; that is to say: it takes reality along only in mirror images of the real, but the real itself it does not touch with a finger. It lets the real world go its own crooked way and builds itself, by way of compensation, an ideal world, richer, more glorious and more enchanting than the one formed of the coarse stuff. And men are delighted by the sight of the newly created world; light and shadow, how happily they are distributed, the refractions of light manifold and the mingling of colors melting, lovely, ravishing. Here nature lays bare her hidden beauties, her virginal secrets. Here one goes at speed over hill and dale; even adversity is a spice, the delays, belles horreurs, are interesting. And the human beings in this world, how open-hearted they are! They confide to us everything they feel and think. One comes so close to them that one can sense their breathing and feel where the pulse beats. Here is advancement, here is experience, here knowledge of men is gained: one has ``Auntie's spectacles'' on, and sees each individual right into the heart.

With the ethical imagination things are far humbler. Inasmuch as imagination takes service with the will, reality must indeed also be apprehended from the side of possibilities; but with the ethical possibilities matters stand in quite another way than with the aesthetic. The only possibilities the will heeds are those that can really be realized. The will looks with the eye of imagination at possibility, but, mark well, for the sake of reality; the ethical will goes from reality
% --- p. 90 ---
\setcounter{footnote}{0}%
\opage{90}through possibility over to a new reality, for it reworks reality.

Is the new reality, then, worked over by the will, more beautiful and richer than the one with which we began? We answer: Yes! only that the beauty is not essentially in the outward but in the inward. If the decisive emphasis is laid on the beautification of the outward, then the aesthetic takes the power from the ethical, and the whole endeavor, even if one labors in the sweat of one's face, takes on a predominant stamp of merely aesthetic imagination.

Here is a deception of imagination to which the world of culture, with its restless labor for the conveniences of life, comfort in the lesser and grand beautifications in the greater, great grounds, great buildings, great monuments, etc., is constantly exposed. The deception by no means consists in this, that work is done in the material sphere for progress and beautification, far from it! The deception is that this labor in an outward direction is set up as an unfailing sign of the ethical. The deception is that the ethical, far from keeping pace with culture, is on the contrary weakened, as the luster of the phenomena of splendor sets the standard and puts the invisible beauty of the good in the shade.

The beauty of the good is, when we lay weight on the contrast to the aesthetic, an invisible beauty. This invisible beauty, to be sure, also has its phenomenon, and can be seen in imagination; but the aesthetic contemplation gives only a passing mood: true insight is conditioned by personal experience. Faithfulness, steadfastness, patience, long-suffering are invisible beauties. It is not the outward surroundings that are to shine here; on the contrary! even the aesthetic
% --- p. 91 ---
\setcounter{footnote}{0}%
\opage{91}imagination understands that they --- for the sake of effect --- had best be as humble as possible. A dark prison, a poor room, a sickbed, these are surroundings that suit the invisible beauty.

The aesthetic imagination can poetize patience and fall in love with its poem; but the essence of real patience it does not grasp. It can fall in love with a beautiful face; but when it falls in love with the beauty of patience, a veil falls over its eyes, for it sees not the essence itself but only a semblance, and does not comprehend how the patient soul can in reality endure being so patient. This question only the ethical, indeed strictly speaking the religious-ethical, imagination can understand, but --- on the condition of experience.

We began by looking about us in the great realm of imagination, and now that we have seen that this realm is infinite, our short wandering is at an end. Should anyone say to himself: it is a pity, for now it was just worth beginning, then it would be to me not merely a reminder that this presentation certainly has its imperfection, but also a testimony that on our short way we have nevertheless met a good many things with which it would well be worth making closer acquaintance.

\end{document}
